\documentclass[11pt,letterpaper]{article}
\usepackage[T1]{fontenc}
\usepackage{lmodern}
\usepackage[margin=1in,headheight=15pt]{geometry}
\usepackage{amsmath,amssymb,amsthm,mathtools}
\usepackage{microtype,booktabs,longtable,array,enumitem,xcolor,listings}
\usepackage{xurl,hyperref,fancyhdr}
\definecolor{ink}{HTML}{15354B}
\definecolor{accent}{HTML}{126E82}
\definecolor{papergray}{HTML}{F1F5F7}
\hypersetup{colorlinks=true,linkcolor=ink,citecolor=accent,urlcolor=accent,
 pdftitle={Continuous Probability in Intuitionistic Primitive Recursive Arithmetic},
 pdfauthor={Research prepared for Vladimir Reshetnikov}}
\pagestyle{fancy}
\fancyhf{}
\fancyhead[L]{\small\color{ink}Continuous probability in iPRA}
\fancyhead[R]{\small Research report}
\fancyfoot[C]{\thepage}
\renewcommand{\headrulewidth}{0.3pt}
\setlength{\parskip}{.35em}
\setlength{\emergencystretch}{2em}
\setcounter{tocdepth}{1}
\setlist{itemsep=2pt,topsep=5pt}
\newtheorem{theorem}{Theorem}[section]
\newtheorem{lemma}[theorem]{Lemma}
\newtheorem{proposition}[theorem]{Proposition}
\newtheorem{corollary}[theorem]{Corollary}
\theoremstyle{definition}
\newtheorem{definition}[theorem]{Definition}
\theoremstyle{remark}
\newtheorem{remark}[theorem]{Remark}
\lstset{basicstyle=\small\ttfamily,columns=fullflexible,keepspaces=true,
 breaklines=true,frame=single,rulecolor=\color{ink!25},
 backgroundcolor=\color{papergray},showstringspaces=false}

\begin{document}
\begin{titlepage}
\color{ink}
{\large\sffamily FOUNDATIONS AND FORMALIZATION}\par
\vspace{1.0cm}
{\Huge\bfseries Continuous Probability\par}
\vspace{.25cm}
{\LARGE in Intuitionistic Primitive\par Recursive Arithmetic\par}
\vspace{.65cm}
{\large Finite rational observations, coherent random variables,\par
 certified integration, and quantitative limit theorems\par}
\vspace{.8cm}
\color{black}
\begin{abstract}
We construct a foundation for continuous probability whose finite content
uses only primitive recursive programs and quantifier-free induction.
One coherent algebra of fair-bit cylinder tables supports represented
events, arbitrary random variables in the bounded convergence-in-probability
metric, and stronger $L^1$ and $L^2$ representations. Error bounds are part
of the data of each limiting construction.

The development includes countable additivity with explicit tails,
integration and quantitative Vitali convergence, independent copies,
conditional averaging and Bayes updating, and continuous-state transitions.
Uniform, exponential, normal, Cauchy, and singular Cantor laws are treated.
A dyadic histogram construction realizes effective densities coherently on
the common bit space while preserving certified moment tails. The Gaussian
normalization and moments are proved by finite quadrature estimates.
Explicit weak- and strong-law bounds and an elementary quantitative central
limit theorem are derived, with careful treatment of discontinuous threshold
events and almost-sure exceptional sets.

The archive supplies the complete human argument, the arithmetic coding and
Lean completion plan, an independently compiled finite Lean kernel with
30 named theorems, and exact-rational reference computations. The continuous
theorems are not claimed to be completed embedded iPRA proof objects; the
remaining derivations and the existing repository providers are identified.
\end{abstract}
\vfill
{\large Prepared for Vladimir Reshetnikov\par}
{20 November 2026\par}
\vspace{.3cm}
{\small Source snapshots and verification receipts are recorded in the archive.}
\end{titlepage}

\tableofcontents
\newpage
\section{Scope, logical foundation, and the meaning of a theorem}
\label{sec:scope}

The aim is a useful foundation for continuous probability inside
intuitionistic primitive recursive arithmetic, abbreviated iPRA. The main
outcome is an explicit arithmetic interpretation of a substantial represented
fragment: continuous and singular continuous distributions, integration,
products, conditioning with separation data, quantitative convergence, and
laws of large numbers and a central limit theorem. It is not an axiom asserting
that the full classical universe of measurable sets exists inside a
first-order number theory.

The subject outline in the requested probability-theory article motivates
the choice of topics \cite{wikiProbability}. Constructive probability has a
large existing literature; Chan's systematic treatment, for example, develops
probability and stochastic processes in Bishop-style constructive analysis
\cite{chan}. Our additional restriction is proof-theoretic: each computation
must have a primitive recursive description, and the induction used in the
object theory must be quantifier-free. A constructive argument in a stronger
analysis is not automatically an iPRA argument. The following presentation
makes the finite calculations and the quantitative hypotheses visible.

\subsection{The exact object theory}

Use first-order intuitionistic natural deduction with equality, the
connectives $\bot,\wedge,\vee,\to$, and the quantifiers $\forall,\exists$.
Negation abbreviates implication to $\bot$. The nonlogical language has
$0,S$ and a function symbol for each \emph{fixed primitive recursive
definition}. Its axioms are the defining equations, $Sx\ne0$, injectivity
of $S$, and quantifier-free induction:
\[
 A(0,\mathbf u)\wedge
 \forall n\bigl(A(n,\mathbf u)\to A(Sn,\mathbf u)\bigr)
 \ \longrightarrow\ \forall n\,A(n,\mathbf u).
\]
Here $A$ is quantifier-free. A recursion symbol has equations
\[
 f(0,\mathbf u)=g(\mathbf u),\qquad
 f(Sn,\mathbf u)=h(n,f(n,\mathbf u),\mathbf u).
\]
The logical inference rules are the usual intuitionistic introduction and
elimination rules, equality substitution, and weakening/substitution of
derivations. No excluded-middle axiom, Markov principle, countable choice,
or comprehension principle is added. Decidable finite rational comparisons
are justified arithmetic case splits.

This is the theory implemented by ProveIt's
\texttt{ConciseIPRA/PR.lean} and its natural-deduction library
\cite{repoCore}. A typed Lean proof of a theorem about the standard natural
numbers is a different object from a term of that embedded derivation type.
The report keeps these two verification levels separate.

\subsection{Numerical coding and total finite operations}

A signed integer is represented by a pair $(u,v)$ with value $u-v$;
cross multiplication gives equality and order without a quotient operation.
A rational code is $(u,v,d)$ with denominator $d+1$, denoting
$(u-v)/(d+1)$. Redundant codes are harmless. Addition, multiplication,
absolute value, comparison, min/max, and division by a denominator with a
supplied positive lower bound are fixed PR operations on these codes.
For malformed optional data use a declared default value; the validity
predicate guards each theorem. This keeps the underlying operations total.

Pairs and finite lists have standard bounded-search codes. Concretely a
finite list of natural entries can be stored as $(n,B,c)$ with $B\ge2$,
$c=\sum_{i<n}a_iB^i$, and $a_i<B$. Entry extraction is
$\lfloor c/B^i\rfloor\bmod B$. For a freshly constructed list take $B$
larger than its finite maximum. Length, concatenation, sorting by bounded
insertion, prefix sums, products, and table refinement are PR. Exponentially
large tables remain primitive recursive; no claim of computational
efficiency follows just from this classification.

Every bounded Boolean condition has an explicit characteristic function.
For example, a universal row condition becomes a fold of conjunctions over
the coded row count. Its connection with the displayed bounded formula is
proved by quantifier-free induction on that fold. A proof must establish this
reflection statement; running a checker on examples is insufficient.

\subsection{How continuous notation is translated}

The phrase ``let $X$ be a represented random variable'' fixes PR programs for
its approximants, their finite levels, and the required error schedules,
together with proofs of their stated contracts. Natural parameters can remain
free. The statement is a schema over these fixed definitions, or a theorem
over a separately specified bounded program interpreter with a proved clock.
It is not quantification in iPRA over arbitrary functions $\mathbb N\to\mathbb N$.

For example, the regularity of an event name $A_k$ with levels $\ell(k)$ is
the first-order formula
\[
 \forall k,l\quad
 2^{k+l}\operatorname{DiffCount}(A_k,A_l)
 \le 2^{\max(\ell(k),\ell(l))}(2^k+2^l).
\]
The count compares the two tables on their common refinement. Its matrix
is an integer inequality. A finite proof of an output estimate asks for
finitely many instances of the input contracts, uses finite arithmetic, and
then universally generalizes the free precision. One must not perform
induction on the all-precision assertion itself.

Here is an explicit denominator-cleared probability theorem. Let $w_i,a_i$
be natural weights and nonnegative numerators, $W=\sum_{i<n}w_i>0$, and
$D,b>0$. Markov's inequality for values $a_i/D$ and threshold $c/b>0$ is
the single finite inequality
\begin{equation}\label{eq:markov-object}
 cD\sum_{i<n}w_i[ba_i\ge cD]
 \le b\sum_{i<n}w_i a_i.
\end{equation}
The partial-sum version of this same inequality is a quantifier-free
induction invariant. The displayed indicator is a PR Boolean comparison.
After proof, its usual probability interpretation divides by $bDW$.

Similarly, for signed $a_i$, a center $m/(WD)$, and threshold $c/b$,
Chebyshev's power-Markov kernel is
\[
 (cWD)^2\sum_{i<n}w_i[b|Wa_i-m|\ge cWD]
 \le b^2\sum_{i<n}w_i|Wa_i-m|^2.
\]
Only when $m=\sum_iw_ia_i$ has been established is the center the mean.
The continuous results preserve this discipline: all analytic error terms
are replaced by rational allowances before the final arithmetic step.

\subsection{The relation to existing ProveIt work}

Two source snapshots were inspected. The local arithmetic and probability
providers are pinned to
\texttt{b2341af5370cd8a43f72ebb2341c0d7d7395b634}. The newer measure and
observation reports were fetched at the then-current main commit
\texttt{83421593771f877bc1ae98f4229eb35507bff05b}.
The repository's directory README distinguishes the checked
\texttt{ConciseIPRA}, \texttt{Derivations}, and \texttt{Bridge} libraries
from research packets retained as delivered. Several newer packets report
their own checked finite Lean kernels, which still do not amount to full
embedded analytic derivations. No repository-wide rebuild or audit is claimed.

This work reuses the coherent-partition and completion approach, excess-tail
clipping, and finite observation discipline from those newer reports
\cite{repoObservable,repoMeasureTransport}. The existing expectation,
normal-distribution, and CLT packets provide useful thematic context
\cite{repoProbability}. Their finite algebra is partly implemented by actual
\texttt{Prf PRAAx} providers. The new package connects these ingredients into
a continuous probability foundation, gives an explicit coherent density
realization, and supplies a self-contained conservative CLT proof. It makes
no claim to have invented the underlying classical or constructive theorems.

\section{The finite analysis needed for probability}
\label{sec:finite-analysis}

The analytic input is small and quantitative. This section specifies the
finite content of integration and calculus used later, so that a reference
to an integral does not silently assume a classical measure construction.

\subsection{Compact quadrature as a Cauchy construction}

A continuous-function presentation on a rational interval $[a,b]$ supplies
a rational evaluator $\widehat f(x,p)$ at rational $x$, with error at most
$2^{-p}$, a rational bound $M$, and a PR modulus $\omega$ satisfying
\[
 |x-y|\le2^{-\omega(k)}\ \Longrightarrow\ |f(x)-f(y)|\le2^{-k}.
\]
These are contracts on regular real names, with rational guards in the
finite implementation. Piecewise presentations carry a finite list of
breakpoints and one-sided values. The integrals of the functions considered
here are constructed only from such data or from explicitly stronger norm
approximations.

For a partition of mesh at most $2^{-\omega(k)}$, tagged rational quadrature
with evaluation error $2^{-p}$ has uncertainty at most
\[
 (b-a)(2^{-k}+2^{-p}).
\]
To prove this, refine two proposed partitions to their finite common list of
endpoints. Each tag differs from any point in its cell by the mesh bound, so
the finite upper and lower rectangle bounds enclose its value with the
displayed error radius. Their symmetric enclosure has twice that radius;
comparing two quadratures adds their error radii. Multiplying by cell lengths
and summing bounds the difference between sufficiently fine quadratures. It produces a
regular real name for the integral. No supremum or infimum over arbitrary
functions is used. A Lipschitz constant $L$ gives the simpler mesh allowance
$L(b-a)h$.

All continuous operations used below come with local computable bounds.
The integral is positive, additive over adjacent intervals, and linear,
because these statements hold for the common finite sums and their errors
can be made arbitrarily small. An improper integral is admitted only when
its finite annular tails have a PR Cauchy bound. In particular the hypothesis
(K1) below uses finite intervals and is not a circular assumption about an
already-existing infinite integral.

\subsection{Finite calculus identities}

The form of the fundamental theorem used here is the following. Suppose
$g,g'$ have the above presentations and a proved local increment estimate
\[
 |g(x+h)-g(x)-hg'(x)|\le Ah^2/2
\]
on a compact interval. Sum it over a partition of mesh $h_0$. The left
increments telescope, so
\[
 \left|g(b)-g(a)-\sum_i g'(x_i)(x_{i+1}-x_i)\right|
 \le A(b-a)h_0/2.
\]
The explicit vanishing error proves $\int_a^b g'=g(b)-g(a)$.
More generally one may supply a proved increment contract
$|g(x+h)-g(x)-hg'(x)|\le |h|\rho(|h|)$ with a PR decay modulus for $\rho$;
the error is then $(b-a)\rho(h_0)$. A bare continuity assertion about a
derivative is not substituted for this quantitative contract. Integration by parts follows by applying
the same finite increment calculation to the product $fg$ and using its
derivative identity. The exponential increment and derivative estimates
follow from its finite power series and geometric remainders; rational and
polynomial functions use exact divided differences and denominator margins.
Thus the applications to normal and exponential kernels require no general
differentiation oracle.

The one-dimensional substitutions used later also have a finite proof.
For a monotone differentiable map $\psi$ on $[a,b]$, carry a bound for
$|\psi'|$, a derivative modulus, and the local modulus of $f$ on its image.
Subdivide the domain. Replace each image-cell length
$\psi(x_{i+1})-\psi(x_i)$ by
$\psi'(x_i)(x_{i+1}-x_i)$ using the increment bound. Oscillation of $f$
on each image cell controls its tagged-sum error. The resulting explicit
error tends to zero and gives
$\int_{\psi(a)}^{\psi(b)}f=\int_a^b(f\circ\psi)\psi'$.
Approximating represented endpoints contributes at most the endpoint error
times a local integrand bound. In this report the maps are affine maps or
reciprocal maps on intervals separated from zero, where every such bound
is a displayed rational expression. For decreasing maps reverse the
endpoints first, or retain oriented interval lengths throughout the finite
sum identity. Improper substitutions separately carry
their spatial-tail estimates.

On compact rectangles, Fubini is equality of finite double sums followed by
the same oscillation bounds. For an unbounded product integral we first
truncate so the \emph{absolute} omitted integral has a supplied bound. This
is the finite Fubini fragment used for products and Gaussian normalization.
It does not make the value of an arbitrary $L^1$ equivalence class on one
specified measure-zero section meaningful.

For the numerical normal example in the archive, the composite midpoint
error is particularly convenient. A twice differentiable function with
$|f''|\le A$ satisfies
\[
 \left|\int_a^b f-h\sum_i f(x_i+h/2)\right|
 \le A(b-a)h^2/24.
\]
Integrate its second-order pointwise remainder symmetrically on each cell;
the linear term cancels, and $\int_{-h/2}^{h/2}t^2/2\,dt=h^3/24$.
For $w(x)=e^{-x^2/2}$,
$|w''(x)|\le (1+x^2)e^{-x^2/2}\le3$, using $te^{-t}\le1$.
This explains the exact enclosure calculation rather than treating numerical
agreement with a floating-point library as evidence of proof.

% Mathematical section. Citation keys repoObservable and repoMeasureTransport
% should identify the currently fetched November 20 ProveIt packets.
\section{A common probability space and its represented completions}
\label{sec:prob-foundation}

The coherent-partition and completion interfaces in the current ProveIt
framework already provide the appropriate starting point
\cite{repoObservable,repoMeasureTransport}. We specialize them to probability,
add an explicit distinction between arbitrary random variables and integrable
ones, and spell out the certificates needed by discontinuous event operations.
These are constructions on one common probability space. A list of laws with
unrelated pairwise couplings would not supply the joint structure required by
addition, independence, conditioning, or events involving several variables.

\subsection{The exact finite base}

For $n\in\mathbb N$, let $\Omega_n=\{0,1\}^n$, coded by the integers below
$2^n$. Each word has mass $2^{-n}$. A cylinder event is a Boolean table on
$\Omega_n$; a rational cylinder variable is a rational table on $\Omega_n$.
When $m\ge n$, lift a table to $\Omega_m$ by ignoring the last $m-n$ bits.
Every row is then repeated $2^{m-n}$ times. Two or more tables are compared
after lifting them to the largest of their levels. The projection from level
$l$ to $n$ is the composite of the projections through any intermediate
level $m$, so these joint refinements are coherent.

For finite tables write
\[
 \mathbb P_n(A)=2^{-n}\sum_{\omega\in\Omega_n}\mathbf1_A(\omega),
 \qquad I_n(f)=2^{-n}\sum_{\omega\in\Omega_n}f(\omega).
\]
Repetition preserves both quantities. Thus the subscript may be omitted once
a common level is understood. Every construction here is primitive recursive:
the number of rows may be very large, but there is an explicit finite bound.
Rationals are signed integers over positive denominators; comparison is exact
cross multiplication. The denominator guard is retained in every formula.

\begin{lemma}[The finite probability kernel]\label{lem:prob-finite}
The cylinder algebra has exact finite additivity, monotonicity, complement
and union laws. Rational integration is positive and linear, and
\[
 |I(f)-I(g)|\le I(|f-g|),\qquad
 I(|f+g|)\le I(|f|)+I(|g|).
\]
If $\varphi$ is rationally evaluated and $L$-Lipschitz, then
$I(|\varphi(f)-\varphi(g)|)\le L I(|f-g|)$.
For finite tables on independent blocks of bits,
\[
 I(f(\omega)g(\eta))=I(f)I(g).
\]
\end{lemma}
\begin{proof}
The event statements follow from Boolean identities on each row. For example,
$\mathbf1_{A\cup B}+\mathbf1_{A\cap B}=\mathbf1_A+\mathbf1_B$.
The integral inequalities follow by multiplying the corresponding rational
pointwise inequalities by the nonnegative row mass and summing. For the last
identity, write the expectation as a finite double sum and use distributivity:
$\sum_{i,j}f_i g_j=(\sum_i f_i)(\sum_j g_j)$. Normalizing by the product of
the two word counts gives the assertion.

For the iPRA proof, clear the positive denominators and use the partial-sum
inequality as the induction invariant. It is quantifier-free. Row conditions
are decidable Boolean tests; their bounded conjunction is encoded by a Boolean
fold. Universal presentation hypotheses are instantiated at the finitely many
rows used. No induction on a statement about all precisions is needed.
\end{proof}

Natural-weight finite probability spaces are also available: replace
$2^{-n}$ by $w_i/W$, where $W=\sum_iw_i>0$. They are useful for rational
quadrature laws. Their mass vectors need not be exactly dyadic, so they should
not be silently identified with a single finite fair-bit table. They can be
approximated by dyadic masses with explicit rounding error, or retained as
finite auxiliary laws with separately supplied transports.

\subsection{Three different distances}

On rational cylinder variables define
\[
 d_0(f,g)=I(\min(1,|f-g|)),\qquad
 d_1(f,g)=I(|f-g|),\qquad
 d_2(f,g)^2=I((f-g)^2).
\]
The first two are exactly rational. Square roots in $d_2$ may instead be
handled by squared rational inequalities, or by rational bisection with a
proved error bound. On cylinder events define
\[
 d_{\rm ev}(A,B)=\mathbb P(A\mathbin{\triangle}B)
                =d_1(\mathbf1_A,\mathbf1_B).
\]
The triangle inequality for $d_0$ follows from
\[
 \min(1,|x-z|)\le\min(1,|x-y|)+\min(1,|y-z|).
\]
The triangle inequality for $d_1$ is Lemma~\ref{lem:prob-finite}. The usual
$d_2$ inequality follows from finite Cauchy--Schwarz; its square-free kernel is
\[
 \left(\sum_i w_i a_i b_i\right)^2
 \le\left(\sum_iw_i a_i^2\right)\left(\sum_iw_i b_i^2\right),
\]
proved by expanding the nonnegative sum
$\sum_{i,j}w_iw_j(a_i b_j-a_j b_i)^2$.

\begin{lemma}[Finite probability interpretation of $d_0$]
For rational $0<\delta\le1$,
\[
 \delta\,\mathbb P(|f-g|>\delta)
 \le d_0(f,g)
 \le \delta+\mathbb P(|f-g|>\delta).
\]
\end{lemma}
\begin{proof}
On a row with $|f-g|>\delta$, the capped difference is at least $\delta$;
on every other row it is at most $\delta$. On exceptional rows it is at most
one. Sum these two elementary inequalities.
\end{proof}

Thus $d_0$ controls convergence in probability without requiring the first
absolute moment. This distinction is essential for continuous laws such as
the Cauchy distribution. An $L^1$-only foundation would exclude its identity
random variable even though its distribution and all bounded observables are
well defined.

\subsection{Names, equality, and their arithmetic meaning}

Put $\varepsilon_k=2^{-k}$. For any one of the displayed metrics, a
\emph{regular name} is a fixed primitive-recursive family $x_k$ of finite
tables, with their levels, satisfying
\begin{equation}\label{eq:prob-regular}
 d(x_k,x_l)\le\varepsilon_k+\varepsilon_l\qquad(k,l\in\mathbb N).
\end{equation}
Its limit notation denotes a represented object. Two names are equal when
\begin{equation}\label{eq:prob-equal}
 d(x_k,y_l)\le\varepsilon_k+\varepsilon_l\qquad(k,l\in\mathbb N).
\end{equation}
This is an arithmetical all-precision relation, not a Boolean equality test.
It needs no quotient type inside iPRA.

For completeness, the small rational fact used by this definition is worth
recording. If rational $a,b$ satisfy $a\le b+2^{-m}$ for all $m$, then $a\le b$.
Decide their rational comparison. If $a>b$ and $D$ is a denominator of their
difference, then $a-b\ge1/D$; choosing $m=D+1$ contradicts the displayed
inequality because $2^{D+1}>D$. This is a primitive-recursively bounded
rational argument. Applying it to a triangle inequality with the intermediate
name evaluated at precision $m$ proves transitivity of
\eqref{eq:prob-equal}. Reflexivity and symmetry are immediate.

An $L^0$ random variable means a $d_0$ name. An integrable random variable is
a $d_1$ name; a square-integrable random variable is a $d_2$ name. Their
forgetful maps are justified by $d_0\le d_1\le d_2$, the second inequality
being finite Cauchy--Schwarz with total mass one. An event means a
$d_{\rm ev}$ name of Boolean tables. It is identified modulo null distance;
it does not carry a general procedure deciding whether an individual infinite
bit sequence belongs to it.

These names form the represented part of the usual external metric
completions. They are closed under supplied jointly PR Cauchy families with
explicit PR moduli, by a diagonal precision choice. They are not claimed to
form a classically complete space under arbitrary Cauchy sequences. The
ordinary measure-theoretic interpretation is an adequacy statement; it is
not an additional axiom of the internal arithmetic.

\begin{proposition}[Continuous finite operations extend]
Suppose a fixed PR approximate table operation $F_p$ satisfies
\[
 d(F_p(a),F_q(b))\le Ld(a,b)+\varepsilon_p+\varepsilon_q
\]
for a supplied rational $L\ge0$. If $L\le2^h$, then
$y_k=F_{k+2}(x_{k+h+2})$ is a regular output name, depending only on the
represented input.
\end{proposition}
\begin{proof}
Substituting \eqref{eq:prob-regular} makes the right-hand side at most
$(\varepsilon_k+\varepsilon_l)/2$. Substituting
\eqref{eq:prob-equal} gives the same bound for equivalent inputs. All
precision choices are explicit PR expressions. This is the common completion
rule already used in the repository, now applied to probability tables.
\end{proof}

Consequently addition, rational scaling, absolute value, positive and negative
parts, minimum, maximum, and clipping have their expected representations.
For example addition uses one guard bit because the two errors add. A
bounded Lipschitz test function gives a real expectation from every $L^0$
name: if $|h|\le M$ and $h$ has Lipschitz constant $L$, then
\[
 I(|h(f)-h(g)|)\le \max(L,2M)d_0(f,g).
\]
Indeed, for $|f-g|\le1$ use Lipschitz continuity; otherwise use the bound
$2M$. This provides the law of a general random variable through bounded
observables even when the identity has no expectation.

Every $L^0$ name also supplies effective tightness. Choose one finite
approximation $x_r$ and a rational $B$ bounding all its values. For every
later approximation $x_s$,
\[
 \mathbb P(|x_s|>B+1)\le d_0(x_s,x_r)
                         \le\varepsilon_s+\varepsilon_r.
\]
Thus choosing $r$ from the desired error gives an explicit tail radius.
Effective continuous maps with supplied PR uniform moduli on rational compact
intervals extend to $L^0$: first discard this small tail, then use the local
modulus on the compact part. A reciprocal or logarithm additionally needs
an explicit bound on the probability of approaching its singular set.

\subsection{Integration, events, and the threshold issue}

\begin{theorem}[Represented expectation]
If $X$ has an $L^1$ name $x_k$, then $I(x_k)$ is a regular rational real
name. It defines $\mathbb E X$, independent of the name, and
\[
 |\mathbb E X-\mathbb E Y|\le\|X-Y\|_1.
\]
The resulting expectation is linear and positive. The completed order is
$X\le Y$ when $\|(X-Y)^+\|_1=0$.
\end{theorem}
\begin{proof}
Apply the exact finite inequality
$|I(x_k)-I(x_l)|\le d_1(x_k,x_l)$ and then
\eqref{eq:prob-regular}. Equivalent names give equivalent scalar names by
the same inequality. Every linear identity holds exactly on a common finite
table and passes to names with its explicitly budgeted errors. Positivity
follows because the negative part has zero norm. This uses order on represented
classes, not a decision about pointwise order of completed values.
\end{proof}

\begin{proposition}[Represented Boolean event algebra]
Complement is an isometry for $d_{\rm ev}$, and union and intersection satisfy
\[
 d_{\rm ev}(A\cup B,A'\cup B'),\quad
 d_{\rm ev}(A\cap B,A'\cap B')
 \le d_{\rm ev}(A,A')+d_{\rm ev}(B,B').
\]
They extend to represented events, preserving all Boolean identities modulo
null distance. Probability is $\mathbb P(A)=\mathbb E\mathbf1_A$ and satisfies
the finite probability laws at every precision.
\end{proposition}
\begin{proof}
Changing either Boolean operation can only affect a row where at least one
input indicator changes. The finite union bound proves the estimate.
Complement uses the same level; the binary operations use one guard bit.
Their exact finite identities and the expectation contraction then prove the
asserted completed identities.
\end{proof}

There is an essential qualification. The map $x\mapsto\mathbf1_{x\le t}$
is discontinuous. Small $L^0$ or $L^1$ distance alone does not make its
threshold indicators close. For instance the constant tables $1/n$ converge
to zero, while $\mathbf1_{1/n\le0}$ is always zero and
$\mathbf1_{0\le0}$ is one.

\begin{lemma}[Threshold lifting with a boundary certificate]
Let $x_r$ be an $L^0$ name and let $t$ be rational. Suppose PR functions
$b(k),R(k)$ satisfy, at every $r\ge R(k)$,
\[
 \mathbb P(|x_r-t|\le2^{-b(k)})\le2^{-k}.
\]
Then the threshold events $\{x_r\le t\}$ can be reindexed by a PR function
to form a regular represented event.
\end{lemma}
\begin{proof}
For two finite tables $x_r,x_s$ and $0<\delta\le1$, a change of threshold
indicator implies either $|x_r-t|\le\delta$ or $|x_r-x_s|>\delta$.
Consequently
\[
 d_{\rm ev}(\{x_r\le t\},\{x_s\le t\})
 \le\mathbb P(|x_r-t|\le\delta)+\delta^{-1}d_0(x_r,x_s).
\]
At requested precision $k$, choose $\delta=2^{-b(k+3)}$ and choose a
monotone precision schedule $r(k)$ at least
$R(k+3)$ and $b(k+3)+k+5$. For $l\ge k$ both approximation errors are
at most $2^{-r(k)}$, so the right-hand side is at most
$2^{-(k+3)}+2^{-(k+4)}<2^{-k}$. This is stronger than the required
regularity. A represented real threshold is treated by budgeting its rational
approximation error inside the boundary band.
\end{proof}

A locally bounded density supplies the band estimate through
$\mathbb P(t-\delta\le X\le t+\delta)\le2M\delta$ when its density is
bounded by $M$ there. Its finite sampling or transport errors are included
in the certificate. The same lemma can apply to singular continuous laws;
having a density is sufficient, not necessary. Probabilities of hard events
in later statements always require their event names or a stated band
certificate. Finite rational threshold inequalities themselves need no such
extra assumption.

\subsection{Countable operations retain their tail data}

Let $(A_j)$ be a fixed jointly PR family of represented events and let
$H_N=\bigcup_{j<N}A_j$. Suppose a PR function $C$ supplies
\begin{equation}\label{eq:union-cauchy}
 M\ge N\ge C(k)\quad\Longrightarrow\quad
 d_{\rm ev}(H_M,H_N)\le2^{-k}.
\end{equation}
This is a finite-union contract; it makes no reference to an already
constructed infinite union. The regularized finite unions define a
represented event $\bigcup_jA_j$. To obtain an actual cylinder table at
output precision $k$, take $N=C(k+3)$ after monotone closure, approximate
each of its $N$ events with error at most $2^{-(k+3)}/\max(1,N)$, and take
their finite union. The union estimate and \eqref{eq:union-cauchy} give the
required error budget. The schedule is PR; the coarser bound of $N$ additional
guard bits is sufficient if one wishes to avoid logarithms.

A useful sufficient hypothesis is
\[
 M\ge N\ge C(k)\quad\Longrightarrow\quad
 \sum_{j=N}^{M-1}\mathbb P(A_j)\le2^{-k},
\]
by the finite union bound. If the $A_j$ are pairwise disjoint, their finite
union probabilities are the corresponding partial sums; passing through the
same rate proves countable additivity for this presentation. Intersections
follow by complements. This is a substantive countable theorem, but it does
not assert that every countable sequence automatically comes with a tail
modulus.

\subsection{A quantitative convergence theorem without hidden thresholds}

Define $T_K(x)=\max(-K,\min(K,x))$ and
$R_K(x)=(|x|-K)^+$ for rational $K\ge0$. These are continuous, 1-Lipschitz
functions and $|x-T_K(x)|=R_K(x)$.

\begin{lemma}[The clipping estimate]\label{lem:prob-clipping}
For finite rational tables,
\[
 d_1(f,g)\le\max(1,2K)d_0(f,g)+I(R_K(f))+I(R_K(g)).
\]
\end{lemma}
\begin{proof}
Insert $T_K(f)$ and $T_K(g)$ between $f$ and $g$ and use the triangle
inequality. Since clipping is nonexpansive and bounded by $K$,
\[
 |T_K(x)-T_K(y)|\le\min(|x-y|,2K)
                    \le\max(1,2K)\min(1,|x-y|).
\]
Integrate the resulting pointwise inequality. There are no hard completed
threshold events in this proof.
\end{proof}

\begin{theorem}[Effective Vitali completion]
Suppose $f_n$ is a jointly PR family of finite tables with a PR $d_0$ Cauchy
modulus. Suppose also that a PR threshold $K(k)$ satisfies
$I(R_{K(k)}(f_n))\le2^{-k}$ for every $n$. Then the same family has a PR
$d_1$ Cauchy modulus and hence an integrable represented limit.
\end{theorem}
\begin{proof}
For desired error $2^{-k}$ choose $K=K(k+2)$, and choose $h$ with
$\max(1,2K)\le2^h$. Request $d_0(f_m,f_n)\le2^{-(k+h+2)}$ from the
given modulus. Lemma~\ref{lem:prob-clipping} gives
$d_1(f_m,f_n)\le3\cdot2^{-(k+2)}<2^{-k}$. This computes the claimed
modulus, with no extraction from qualitative convergence.
\end{proof}

The same proof gives convergence of an already represented $L^1$ family to
an $L^0$ limit under uniform excess-tail bounds. An $L^1$ majorant supplies
such bounds constructively. If $G\ge0$ has finite approximation $g_r$ and
$K$ bounds every value of $g_r$, then
\[
 \mathbb E(G-K)^+\le\|G-g_r\|_1.
\]
Indeed $(G-K)^+\le|G-g_r|$ in the completed lattice. Therefore
$|X_n|,|X|\le G$ gives uniform tail control for all $X_n$ and $X$.
Combining this with a supplied PR convergence-in-probability rate proves a
quantitative dominated-convergence theorem. This specializes the existing
repository measure framework; neither domination alone nor pointwise
convergence alone is being used to manufacture the missing rate.

\subsection{Conditional expectation on the dyadic filtration}

Let $\mathcal F_n$ denote observation of the first $n$ bits. For a table $f$
at level $m\ge n$, define $C_n f$ by averaging its $2^{m-n}$ entries above
each length-$n$ prefix. If $m<n$, simply lift $f$. Every denominator is a
positive power of two.

\begin{theorem}[Conditional averaging and a quantified martingale limit]
The finite maps $C_n$ preserve positivity and expectation, are $d_1$
contractions, and satisfy $C_nC_m=C_{\min(n,m)}$. They therefore extend to
represented $L^1$ variables. If $X$ has regular name $x_k$, with $x_k$
depending on $\ell(k)$ bits, then
\[
 n\ge\ell(k)\quad\Longrightarrow\quad
 \|C_nX-X\|_1\le2\varepsilon_k.
\]
In particular the sequence $C_nX$ is a represented martingale converging to
$X$ with an explicit PR rate.
\end{theorem}
\begin{proof}
The inequality that the absolute value of an average is at most the average
of the absolute values, summed over prefixes, proves contraction. Regrouping
the same finite sums proves preservation of expectation and the tower law.
For an $\mathcal F_n$ cylinder variable $h$, the same finite regrouping gives
$I(hC_n f)=I(hf)$, the defining conditional-expectation identity. Positivity
is termwise. The completion rule extends these facts to names.

If $n\ge\ell(k)$, then $C_n x_k=x_k$. Hence
\[
 \|C_nX-X\|_1
 \le\|C_n(X-x_k)\|_1+\|x_k-X\|_1
 \le2\|X-x_k\|_1\le2\varepsilon_k.
\]
For error $2^{-p}$, the monotone schedule
$N(p)=\max_{j\le p+1}\ell(j)$ suffices. The tower identity says exactly that
$(C_nX)$ is a martingale. This is convergence of the canonical martingale of
a supplied terminal variable, not an unrestricted theorem for arbitrary
bounded martingales.
\end{proof}

Conditioning on a general represented event $B$ is also available when a
rational margin $\beta>0$ with $\mathbb P(B)\ge\beta$ is supplied:
$\mathbb P(A\mid B)=\mathbb P(A\cap B)/\mathbb P(B)$. Approximate the
denominator enough to keep it above $\beta/2$; ordinary rational division
then has an explicit error budget. A null event supplies no such division,
and a general regular conditional distribution is extra data, not a consequence
of this quotient formula.

\subsection{Independent copies, a uniform law, and a singular continuous law}

A primitive-recursive pairing function partitions the countably many bit
positions into disjoint infinite rows. Copy the table for a represented
variable onto row $i$ by reading only positions assigned to $(i,j)$. Each
finite demand touches finitely many positions and is therefore a cylinder
table; counting unused positions preserves its law and its distance bounds.
The finite product identity proves independence of every finite family of
different rows. Thus a supplied represented random variable has a jointly
PR sequence of independent identically distributed copies. This is an explicit
construction, not an appeal to an arbitrary extension theorem.

Let $U_n$ assign the midpoint $(j+1/2)2^{-n}$ to the $j$th dyadic cell.
For $m\ge n$, $|U_m-U_n|\le2^{-n-1}$ on every common row. The resulting
$U$ is a bounded $L^1$ and $L^0$ variable. Counting cells gives
\[
 \left|\mathbb P(U_n\le t)-\max(0,\min(1,t))\right|\le2^{-n}
\]
for rational $t$, so its CDF is the uniform CDF. The same count supplies
threshold boundary bands of mass at most $2\delta+2^{1-n}$. In particular,
singletons have probability zero; this does not make the value of an actual
infinite sample decidable.

For a continuous law that has no density, use the same bits in
\[
 C_n=\sum_{j=1}^n\frac{2b_j}{3^j}.
\]
The remainder is bounded by $3^{-n}$, so $C_n$ is uniformly Cauchy with an
explicit PR rate. Each prefix of length $n$ puts the limit in one of $2^n$
closed intervals of length $3^{-n}$, with mass $2^{-n}$ each. A singleton
lies in at most one such interval and therefore has probability at most
$2^{-n}$ for every $n$, hence zero. More quantitatively, a sufficiently short
band of width below the level-$n$ gaps meets at most two intervals, giving
a boundary-band modulus and hence a represented continuous CDF.

The same $2^n$ intervals cover the support with total ordinary length
$(2/3)^n$. Choosing $n=3k$ makes this at most $2^{-k}$, since
$(2/3)^3<1/2$. Thus the support is Lebesgue-null with an explicit covering
rate, while the constructed law gives it full mass in the sense that every
finite cover stage has probability one. In standard semantics this is the
Cantor singular continuous distribution. It verifies that the foundation
handles continuous laws beyond absolutely continuous density models.

\subsection{What the arithmetic claim does and does not say}

For each fixed presentation, every displayed finite estimate is an arithmetic
formula over natural codes. The proof recipe consists of exact rational
identities, proved Boolean tests, and quantifier-free induction on finite
partial sums, refinements, and loops. The regularity hypotheses are used at
the finitely many indices demanded by the output precision. A theorem about
all precisions is obtained by universal introduction after those finite
proofs; it is not used as an induction matrix. Completed objects are a
convenient readable notation for this family of contracts.

There is no total universal PR function evaluating every arbitrary PR
program code: if $V(e,n)$ were such a function, then $V(n,n)+1$ would be PR
and its own code would give a contradiction. Accordingly, names here are
fixed PR definitions, possibly with numerical parameters. A variable program
code requires a bounded interpreter and a supplied clock with a proof of
sufficiency. The mathematical claims above are proof-construction schemata
with explicit finite targets; a Lean proof of those targets in ordinary
arithmetic is not yet a term of the embedded object-theory type
$\mathrm{Prf}\;\mathrm{PRAAx}$.

Two elementary obstructions explain why the representation data cannot be
discarded. For a machine index $e$, the stage-$n$ assertion ``$e$ has halted
by stage $n$'' is a PR Boolean sequence, monotone and bounded by one. A
uniform procedure producing its regular real limit to error $1/4$ would
decide halting. Thus bounded monotonicity does not uniformly yield a PR
convergence modulus. Similarly, enumerate no open interval until that machine
halts, and then enumerate $(0,1)$. The resulting effective open set has
uniform probability either zero or one. A uniform regular probability for
all such descriptions would again decide halting. Therefore this report does
not identify all effective Borel descriptions with regular event names.

Exact real equality, arbitrary sigma-algebra completion, automatic rates from
qualitative almost-sure convergence, and general disintegration are not hidden
axioms of the construction. Positive separation, tail bounds, boundary-band
bounds, and convergence rates are retained precisely where the finite proof
uses them. This is what keeps the continuous mathematics within the stated
intuitionistic primitive-recursive proof discipline.

% Continuous probability: mathematical section prepared for inclusion in the main report.
% Requires amsmath, amssymb, amsthm. All analytic statements below are presented
% with explicit approximation contracts, not as already compiled Lean theorems.

\section{Continuous distributions with explicit approximation contracts}
\label{sec:continuous-laws}

A law and an integrable random variable are different objects.  A probability
law may have finite total mass and no first moment.  Accordingly, the density
construction below first produces probabilities of bounded regular events and
expectations of bounded test functions.  An expectation of an unbounded
observable is admitted only after a separate absolute-tail certificate is
provided.  This distinction includes the Cauchy law without incorrectly
assigning it an ordinary expectation.

\subsection{From a nonnegative kernel to a probability law}

A \emph{kernel presentation} consists of the following data.  The kernel
$w\colon\mathbb R\to[0,\infty)$ is given locally by primitive recursive rational
approximations, a modulus of continuity on each rational compact interval, and
a rational bound there.  Finitely many explicitly listed discontinuities are
also allowed, by splitting the interval at them.  There is a primitive
recursive radius $R(k)\geq 1$ such that, for every rational $B\geq R(k)$,
\begin{equation}
 \int_{[-B,B]\setminus[-R(k),R(k)]}w(x)\,dx\leq 2^{-k}.
 \tag{K1}\label{eq:kernel-tail}
\end{equation}
All integrals in this hypothesis are finite-interval integrals, computed by
rational quadrature with the supplied modulus.  Finally there are explicit
rationals $z_0>0$ and $R_0>0$ with
\begin{equation}
 z_0\leq\int_{-R_0}^{R_0}w(x)\,dx.
 \tag{K2}\label{eq:kernel-positive}
\end{equation}
The inequalities are finite rational approximation contracts, including any
rounding allowance; neither an unbounded search for a cutoff nor a decision
about equality of real numbers is part of the data.

\paragraph{Normalization theorem.}
These data determine a fast Cauchy name for
$Z=\int_{\mathbb R}w$, with $Z\geq z_0$, and a probability law with density
$q=w/Z$.  If $g$ is bounded by a known rational $C$ and is locally effectively
integrable against $w$, its expectation has a primitive recursive Cauchy name.
For an unbounded $g$ the same conclusion follows from a supplied primitive
recursive cutoff $R_g(k)$ for the tails of $|g|w$.

\begin{proof}
Choose a rational quadrature $a_k$ for the integral over
$[-R(k+2),R(k+2)]$ with error at most $2^{-k-2}$.  Replace $R$ by its finite
maximum up to its argument if necessary, so that it is nondecreasing.
For $\ell\geq k$, compare the two finite integrals on the larger interval.
Their difference is at most $2^{-k-2}$ by \eqref{eq:kernel-tail}.  Thus
$|a_\ell-a_k|\leq 3\cdot2^{-k-2}$.  Reindexing gives the stipulated fast
Cauchy form.  Positivity follows by retaining the fixed interval in
\eqref{eq:kernel-positive}.  Division is effective because its denominator
is separated from zero: for positive $u,v\geq z_0/2$,
\[
 |u^{-1}-v^{-1}|\leq 4z_0^{-2}|u-v|.
\]
The mass outside $[-R(k),R(k)]$ is at most $2^{-k}/z_0$.
For bounded $g$, the discarded expectation is at most $C2^{-k}/z_0$;
for unbounded $g$ use the specified $|g|w$ tail.  All remaining quantities
are finite sums.  Normalization, additivity, and positivity follow by
comparing those sums at a common requested precision.  No assertion that
every measurable function is integrable has been made.
\end{proof}

A convenient additional certificate is a local bound $q\leq K_R$ on
$[-R,R]$.  It controls threshold boundaries:
\[
 \mathbb P\{t-\eta<X\leq t+\eta\}\leq 2K_R\eta
 \qquad (|t|+\eta\leq R).
\]
Therefore bounded intervals have effective event names, and the distribution
function $F(t)=\mathbb P\{X\leq t\}$ is a represented continuous function.
If $q\leq K$ globally, $|F(s)-F(t)|\leq K|s-t|$.  These are the certificates
needed to turn continuous-test statements into assertions about exact
threshold events.

\subsection{An elementary primitive recursive exponential}

For completeness, the exponential evaluations required below can be reduced
to finite rational arithmetic.  On $0\leq t\leq T$, with $T\geq1$ an integer,
use $E_N(t)=\sum_{j=0}^{N}t^j/j!$.  After index $2T$, successive terms have
ratio at most $1/2$.  Also
$T^{2T}\leq 2^{2T(T+1)}$.  Consequently the choice
\[
 N=2T+k+2T(T+1)+2
\]
makes the remaining series tail at most $2^{-k-2}$.  The case $T=0$ is exact.
The reciprocal $1/E_N(t)$ approximates $e^{-t}$ with no greater absolute
error, because $E_N(t)\geq1$.  Multiplication, addition, and the derivative
identities for the resulting names follow from finite series identities and
these tail estimates.  The local bounds used in quadrature can thus be
proved inside the same finite arithmetic development.

The familiar estimates $e^t\geq1+t$ for $t\geq0$, $2\leq e\leq3$, and
$e^{1/2}\leq2$ have direct positive-series proofs.  They give explicit
rational estimates below.  The functions need not be introduced as an
uninterpreted analytic oracle.

\subsection{Uniform and exponential laws}

For rational $a<b$, the uniform density is
\[
 q(x)=\frac{\mathbf1_{[a,b]}(x)}{b-a}.
\]
Splitting at $a,b$ makes quadrature elementary.  Direct integration yields
\[
 \mathbb E X=\frac{a+b}{2},\qquad
 \operatorname{Var}(X)=\frac{(b-a)^2}{12},\qquad
 \mathbb E X^m=\frac{b^{m+1}-a^{m+1}}{(m+1)(b-a)}.
\]
These formulas are finite polynomial identities after denominators are
cleared.  Endpoints are null, so the choice of endpoint convention does not
affect the law.

Let $\lambda>0$ be rational, or a represented parameter with a supplied
rational lower bound $\lambda_0>0$.  The exponential density is
\[
 q_\lambda(x)=\lambda e^{-\lambda x}\mathbf1_{[0,\infty)}(x).
\]
The finite-interval fundamental theorem gives
$\int_0^R q_\lambda=1-e^{-\lambda R}$.  Since $e^{-s}\leq2^{-s}$ for integer
$s\geq0$, a cutoff $R\geq k/\lambda_0$ gives mass tail at most $2^{-k}$.
The derivative on the positive half-line is locally bounded using any
supplied upper bound on $\lambda$, and the jump at zero is a listed
breakpoint.

Integration by parts on $[0,R]$, followed by the explicit boundary and tail
bounds, gives
\begin{equation}
 \mathbb E X^m=\frac{m!}{\lambda^m},\qquad
 \mathbb E X=\lambda^{-1},\qquad
 \operatorname{Var}(X)=\lambda^{-2}.
 \label{eq:exponential-moments}
\end{equation}
Here is a uniform moment-tail certificate, so that this passage is justified
for every fixed natural $m$:
\begin{equation}
 \int_R^\infty x^m q_\lambda(x)\,dx
 \leq 2^{m+1}m!\lambda_0^{-m}e^{-\lambda_0R/2}.
 \label{eq:exponential-tail}
\end{equation}
Indeed $t^m\leq2^m m!e^{t/2}$ follows by retaining the $m$th positive term of
the exponential series.  Substitute $t=\lambda x$ and integrate the remaining
exponential.  For example, if $C$ is any integer at least
$\max(1,2^{m+1}m!\lambda_0^{-m})$, then
$R=2(k+C+1)/\lambda_0$ makes the right side at most $2^{-k}$.
This cutoff is deliberately conservative and is primitive recursive in all
displayed integer and rational inputs.

\subsection{The normal law without a hidden normalization theorem}

Put $w(x)=e^{-x^2/2}$ and define $Z$ by the normalization theorem.  On
$[-1,1]$, $w\geq1/2$, hence $Z\geq1$.  For $R\geq1$,
\begin{equation}
 \int_{|x|>R}w(x)\,dx
 \leq \frac{2}{R}e^{-R^2/2}.
 \label{eq:gaussian-tail}
\end{equation}
To prove this, first integrate only over $[R,B]$.  Since $x/R\geq1$ there,
\[
 \int_R^B e^{-x^2/2}\,dx
 \leq R^{-1}\int_R^B x e^{-x^2/2}\,dx
 \leq R^{-1}e^{-R^2/2}.
\]
Reflect and use the same finite inequality at every $B$.  In particular
$Z\leq4$.  A simple radius $R=2(k+2)$ makes the mass tail at most $2^{-k}$;
the sharper bound \eqref{eq:gaussian-tail} is usually much better.

Define the standard normal density by $\varphi=w/Z$.  It is bounded by $1$
and is globally $1$-Lipschitz.  In fact
$|w'(x)|=|x|e^{-x^2/2}\leq1$, since
$e^{x^2/2}\geq1+x^2/2\geq|x|$.
On every $[-R,R]$ it also has an explicit positive lower bound,
\begin{equation}
 \varphi(x)\geq\frac{1}{4\cdot3^{\lceil R^2/2\rceil}}.
 \label{eq:gaussian-slope}
\end{equation}
This is useful for bounded inverse-distribution calculations; it supplies
the separation that a blind bisection search would otherwise lack.

For all integers $m\geq0$ and $R\geq1$,
\begin{equation}
 \int_{|x|>R}|x|^m\varphi(x)\,dx
 \leq C_m e^{-R/4},\qquad C_m=8\cdot4^m m!.
 \label{eq:normal-moment-tail}
\end{equation}
Indeed $|x|^m\leq|x|^{2m}\leq4^m m!e^{x^2/4}$ on $|x|\geq1$;
then $e^{-x^2/4}\leq e^{-|x|/4}$ and $Z\geq1$.  Thus
$R_m(k)=4(k+C_m)$ is one primitive recursive moment cutoff, since
$C_m\leq2^{C_m}$.

Symmetry proves all odd moments zero.  Integration by parts on $[-R,R]$
and \eqref{eq:normal-moment-tail} prove
\begin{equation}
 \mathbb E G^{2m}=(2m-1)!!,\quad
 \mathbb E G=0,\quad \mathbb E G^2=1,\quad \mathbb E G^4=3.
 \label{eq:normal-moments}
\end{equation}
The convention is $(-1)!!=1$.  Explicitly, $w'=-xw$ gives the recurrence
$\int x^{j+2}w=(j+1)\int x^jw$ after the two boundary terms vanish with
the stated rate.  In particular,
\begin{equation}
 \mathbb E|G|^3\leq\tfrac12(\mathbb E G^2+\mathbb E G^4)=2,
 \label{eq:normal-third}
\end{equation}
because $2|x|^3\leq x^2+x^4$.

The usual notation $Z=\sqrt{2\pi}$ can also be justified within this
framework.  Define $\pi=4\int_0^1(1+t^2)^{-1}\,dt$.  Reciprocal substitution
shows $\int_0^\infty(1+t^2)^{-1}\,dt=\pi/2$.  In the positive quadrant,
substitute $y=tx$ in the product integral for $Z^2/4$ and then integrate in
$x$:
\[
 \frac{Z^2}{4}
 =\int_0^\infty\int_0^\infty
      x e^{-(1+t^2)x^2/2}\,dx\,dt
 =\int_0^\infty\frac{dt}{1+t^2}=\frac\pi2.
\]
This is not an invocation of an unrestricted improper change-of-variables
theorem.  Perform the substitution first on $a\leq x\leq R$, $0\leq t\leq T$.
In the original quadrant, the omitted strip $x<a$ has integral at most $2a$,
the omitted tail $x>R$ is bounded by twice the one-sided Gaussian tail, and
on $x\geq a$ the omitted region $y>Tx$ is contained in $y>Ta$.  Choose
$a,R,Ta$ with the Gaussian bounds.  In the transformed integral, the omitted
$t>T$ tail is at most $1/T$; on $t\leq T$ the omissions $x<a$ and $x>R$
are bounded by $Ta^2/2$ and $T e^{-R^2/2}$ respectively.  One simultaneous primitive recursive choice is
\[
 s=k+4,\qquad a=2^{-2s},\qquad T=2^{3s},\qquad R=2(s+1).
\]
Then $Ta=2^s$, $Ta^2/2=2^{-s-1}$, $1/T=2^{-3s}$, and
$T e^{-R^2/2}\leq2^{3s-2(s+1)^2}\leq2^{-s-2}$.
Each of the other three error terms is at most $2^{-s}$.  The sum of all
six is therefore less than $2^{-k}$.  This supplies a single explicit
finite-stage comparison rather than an unrestricted passage to an improper
integral.  Only finite Fubini and one-dimensional substitutions are used.  The normalization-free definition
$\varphi=w/Z$ and every result above are independent of this identification.

For represented $\mu$ and $\sigma\geq\sigma_0>0$, the affine image
$X=\mu+\sigma G$ has density
\[
 \varphi_{\mu,\sigma}(x)=\frac1{\sigma Z}
       \exp\!\left(-\frac{(x-\mu)^2}{2\sigma^2}\right),
\]
mean $\mu$ and variance $\sigma^2$.  The lower bound on $\sigma$ makes
division and the parameter-dependent error estimates effective.

\subsection{Cauchy: a continuous law with no expectation}

For $w(x)=(1+x^2)^{-1}$ the normalization constant again satisfies
$1\leq Z\leq4$, and reciprocal substitution gives $Z=\pi$ as defined above.
For $R\geq1$,
\[
 \int_{|x|>R}w(x)\,dx\leq 2/R.
\]
Thus $R(k)=2^{k+1}$ is a mass cutoff for the normalized density
$q(x)=1/[\pi(1+x^2)]$.  This density is bounded by $1$, is $1$-Lipschitz,
and on $[-R,R]$ is at least $1/[4(1+R^2)]$.

It has no absolute first moment.  On
$[2^j,2^{j+1}]$, $x/(1+x^2)\geq2^{-j-2}$, so the unnormalized first-moment
integral over this interval is at least $1/4$.  Reflect, sum $M$ such
intervals, and divide by $Z\leq4$ to obtain
\[
 \int_{-2^M}^{2^M}|x|q(x)\,dx\geq M/8.
\]
This finite lower bound excludes any finite absolute-moment name.  The
symmetric integrals of $xq(x)$ are zero, but that principal-value statement
is not an expectation.  A realization of the Cauchy law belongs to the
framework's convergence-in-probability class, not its $L^1$ class.


\section{Realizing a continuous law on the common bit space}
\label{sec:realization}

The normalization construction in Section~\ref{sec:continuous-laws} initially
provides an integration rule on the line. To use it as a random variable on
the common probability space, a joint realization is still needed. The
following construction supplies one. It also separates an effective sampler
in probability from a pointwise inverse-CDF algorithm. Strict positivity of
the density is unnecessary for the first construction.

\subsection{A finite monotone-coupling identity}

Let $\mu,\nu$ be finite laws with rational ordered support points and dyadic
weights. On a common sufficiently fine fair-bit table, allocate consecutive
rows to the atoms in increasing order. Denote the resulting nondecreasing
tables by $Q_\mu,Q_\nu$. Each has exactly its prescribed law. The allocations
are compatible with prefix refinement because all cumulative weights have
dyadic denominators.

\begin{lemma}[Finite monotone coupling]\label{lem:monotone-coupling}
If both laws are supported on $[-K,K]$, then
\[
 I|Q_\mu-Q_\nu|
   =\int_{-K}^{K}|F_\mu(t)-F_\nu(t)|\,dt
   \le 2K\sup_t|F_\mu(t)-F_\nu(t)|.
\]
The displayed integral is an exactly rational step-function sum.
\end{lemma}
\begin{proof}
For rational $x,y$ in the interval,
$|x-y|$ is the length of the set of thresholds where
$\mathbf1_{x\le t}$ and $\mathbf1_{y\le t}$ differ. Subdivide the threshold
interval at the finite combined list of support points and sum this identity
over the common fair-bit table. At every threshold, each of the two sets of
rows is an initial segment. Their symmetric difference therefore has mass
$|F_\mu(t)-F_\nu(t)|$. Interchange the two finite sums. The upper bound
multiplies the interval length by the largest integrand value. Sorting,
prefix sums, allocation, and this proof all use bounded rational operations.
\end{proof}

Write $T_K(x)=\max(-K,\min(K,x))$. For laws with larger support, the
same coupling satisfies
\begin{equation}\label{eq:coupling-clipped}
 d_0(Q_\mu,Q_\nu)
 \le 2K\sup_t|F_\mu(t)-F_\nu(t)|
       +\mu(|x|>K)+\nu(|x|>K).
\end{equation}
Indeed replace both tables by their clipped versions. The replacement costs
at most the two indicated exceptional masses in $d_0$; the clipped CDFs
have the same discrepancy on $(-K,K)$ and identical endpoints outside it.
Apply the lemma and $d_0\le d_1$. This estimate, rather than an unbounded
first-moment transport estimate, allows heavy-tailed laws.

\subsection{An explicit histogram sampler}

Suppose $q$ is a normalized kernel as in Section~\ref{sec:continuous-laws},
with local bounds and continuity moduli, and a PR mass-tail cutoff. Put
$\varepsilon_k=2^{-k}$. Choose a monotone rational $K_k\ge1$ such that
\[
 \int_{|x|>K_k}q(x)\,dx\le\varepsilon_k/16,
 \qquad
 \eta_k=\frac{\varepsilon_k}{64(K_k+1)}.
\]
The sequence $\eta_k$ is nonincreasing. It is essential that the following
\emph{quadrature radius} is a separate, possibly much larger number: choose
$H_k\ge K_k$ with mass tail at most $\eta_k/4$. A single radius need not
meet both budgets for a heavy-tailed law.

Partition $[-H_k,H_k]$, including all listed discontinuities, into $J_k$
rational cells of width at most $1$ and with mass at most $\eta_k/4$ each.
A supplied bound $L_{H_k}$ for the density makes any mesh of width at most
$\eta_k/[4\max(1,L_{H_k})]$ sufficient. At each cell $I_j$, quadrature
and downward dyadic rounding give a dyadic number
\[
 0\le p_j\le\int_{I_j}q,
 \qquad
 \sum_j\left(\int_{I_j}q-p_j\right)\le\eta_k/4.
\]
For example, choose a common binary denominator so that each integration
error plus rounding loss is at most $\eta_k/(4\max(1,J_k))$.
Use the cell midpoint as its atom and add an atom at zero of weight
$1-\sum_jp_j$. These are nonnegative dyadic weights summing \emph{exactly}
to one. There is no renormalization of an uncertain total mass.

Let $\mu_k$ be this finite law. Moving mass within one cell changes its CDF
by at most the mass of that cell; moving all omitted and rounded mass to
zero changes it by at most the total omitted mass. Consequently
\begin{equation}\label{eq:histogram-cdf}
 \sup_t|F_{\mu_k}(t)-F(t)|\le3\eta_k/4\le\eta_k.
\end{equation}
The comparison is proved from finite sums and then from their rational
enclosures; the supremum is just notation for the uniform inequality.
Sorted consecutive allocation turns $\mu_k$ into a cylinder table $Q_k$.

\begin{theorem}[Coherent $L^0$ realization of an effective density]
\label{thm:density-realization}
The tables $Q_k$ form a regular $L^0$ name. Their represented law is $q(x)\,dx$.
In particular the normal, exponential, and Cauchy density presentations
produce random variables on the same fair-bit probability space. Independent
copies are available for all three; conditional expectation applies to their
represented integrable observables.
\end{theorem}
\begin{proof}
For $l\ge k$, the two CDFs differ by at most $\eta_k+\eta_l$. Also
\[
 \mu_j(|x|>K_k)\le\varepsilon_k/16+2\eta_j
 \quad(j=k,l),
\]
by the two one-sided CDF comparisons. Apply
\eqref{eq:coupling-clipped} at the \emph{target} radius $K_k$:
\[
 d_0(Q_k,Q_l)
 \le 2(K_k+1)(\eta_k+\eta_l)+\varepsilon_k/8
 \le\varepsilon_k/16+\varepsilon_k/8
 <\varepsilon_k.
\]
Every table is finite and its allocation uses a declared dyadic denominator;
its prefix comparison with another table is therefore the very monotone
coupling used in the proof. This supplies coherence, not merely the existence
of a separate coupling for each pair of laws.

For a bounded Lipschitz test $h$, compare its finite expectation with the
quadrature integral. On matched cell mass the discrepancy is at most its
Lipschitz constant times the mesh width. On omitted and rounded mass it is
at most $2\|h\|_\infty$ times that mass. These errors tend to zero with
explicit PR rates. The expectation stability for $L^0$ tests then identifies
the limit with $\int hq$. This defines the law without a measure-extension
oracle. It also proves equality of the bounded-observable and density
presentations.

For a threshold $t$, the finite band probabilities satisfy
\[
 \mu_k([t-\delta,t+\delta])
 \le\int_{t-\delta}^{t+\delta}q+2\eta_k.
\]
For the left endpoint use the left limit of the finite CDF; continuity of
$F$ and rational slack removal justify the same discrepancy bound there.
A local density bound now gives the boundary certificate from
Section~\ref{sec:prob-foundation}, constructing an event name with probability
$F(t)$. Thus the threshold claim does not precede the sampler construction.
\end{proof}

If a moment-tail certificate for $|x|^p q(x)$ is supplied, this sampler also
has a corresponding integrability refinement. For $p=1,2$ it gives an
$L^p$ name after a PR reindexing. Here are the finite bounds that establish
the assertion rather than merely identify a putative moment. Since the mesh
width is at most one and each assigned cell mass is no larger than its true
mass, for $K>1$,
\[
 \mathbb E_{\mu_j}\bigl(|x|^p\mathbf1_{|x|>K}\bigr)
 \le 2^{p-1}\int_{|x|>K-1}(|x|^p+1)q(x)\,dx.
\]
The zero residual atom contributes nothing. Thus the assigned masses have
uniform effective moment tails. For $p=1$, apply the clipping/Vitali theorem.
For $p=2$, $K\ge1$ and finite tables $u,v$ give
\[
 I|u-v|^2\le12K^2d_0(u,v)+3I(R_K(u)^2)+3I(R_K(v)^2).
\]
Here $R_K(x)=(|x|-K)^+$. Insert the clipped values, use
$(a+b+c)^2\le3(a^2+b^2+c^2)$, and bound their clipped squared difference
by $4K^2\min(1,|u-v|)$. Uniform second-moment tails then give an explicit
$d_2$ Cauchy rate. More generally continuous polynomial observables have
$L^1$ names under their stated absolute-moment tails. The Cauchy identity
variable has no such first-moment refinement, as proved in the distribution
section. For exact moment agreement, first use the bounded Lipschitz test
$x\mapsto(T_Rx)^p$. Its sampler expectation equals its density integral by
the law identification already proved. The uniform moment tails on both
sides permit $R$ to grow with a PR schedule, proving
$\mathbb E Q^p=\int x^p q(x)\,dx$, and likewise for absolute moments.

\subsection{Pointwise quantiles require a different certificate}

A CDF can be continuous while having flat intervals. The preceding sampler
still works, but a unique pointwise inverse need not vary continuously in
its input data. For a safe inverse construction, suppose
\[
 F(b)-F(a)\ge c(b-a)\quad(-R\le a\le b\le R),\qquad c>0,
\]
and the target rational mass $u$ lies between $F(-R)$ and $F(R)$ with
certified margins. To locate its inverse to distance $\epsilon$, evaluate
$F(m)$ at the current midpoint to error $\delta=c\epsilon/4$. If the
enclosure is below $u$, retain the upper half; if it is above $u$, retain
the lower half. If it straddles $u$, its midpoint value has residual at most
$2\delta$ and hence spatial error at most $\epsilon/2$ by the lower-slope
inequality. Otherwise after a prescribed number of halvings the bracket
width is at most $\epsilon$. No comparison decides whether two real
numbers are exactly equal. The normal and Cauchy lower density bounds in
Section~\ref{sec:continuous-laws} supply $c$ on any certified compact range.
Existence is supplied by this construction itself: in the straddling case
intersect the current bracket with $[m-\epsilon/2,m+\epsilon/2]$; the
lower-slope estimate preserves the endpoint signs. Brackets produced at
different precisions cannot be disjoint, since strict increase would
contradict their opposite endpoint inequalities. Their midpoints are
therefore Cauchy with an explicit width bound, and continuity makes their
limit a solution. The same lower-slope estimate proves uniqueness.

\subsection{Moments, covariance, and products of integrable variables}

For $X,Y\in L^2$, the product has an $L^1$ name. In fact finite
Cauchy--Schwarz gives
\[
 \|XY-X'Y'\|_1
 \le\|X\|_2\|Y-Y'\|_2+\|Y'\|_2\|X-X'\|_2.
\]
A bound for the norm of any regular name follows from one of its finite
entries and its error. Choose input precisions with these bounds and the
displayed inequality supplies the product's PR modulus. The same argument
extends Cauchy--Schwarz. Thus
\[
 \operatorname{Var}(X)=\mathbb E(X-\mathbb EX)^2\ge0,
 \qquad
 \operatorname{Cov}(X,Y)=\mathbb E(XY)-\mathbb EX\,\mathbb EY
\]
are represented quantities. Expanding the finite squares proves all familiar
addition identities, and
$|\operatorname{Cov}(X,Y)|^2\le\operatorname{Var}(X)\operatorname{Var}(Y)$.
For every finite rational vector $(a_i)$,
$\sum_{i,j}a_i a_j\operatorname{Cov}(X_i,X_j)
=\mathbb E(\sum_i a_i(X_i-\mathbb EX_i))^2\ge0$.
This gives positive semidefiniteness of covariance matrices without importing
a classical space of random functions. An arbitrary product of two $L^1$
variables is not covered; the $L^2$ or boundedness hypotheses do real work.

There is a second useful closure rule: a represented bounded multiplier times
an $L^1$ variable is integrable, even when the multiplier converges only in
probability rather than uniformly. For finite $|a|,|b|\le M$ and $K\ge0$,
\[
 |ax-by|\le M|x-y|+K|a-b|+2M(|y|-K)^+.
\]
Use the $L^1$ error of the unbounded factor, the bounded multiplier's $L^1$
error, and the automatically supplied excess tail of $y$ to allocate these
three terms. For a bounded $L^0$ multiplier choose approximants clipped to
$[-M,M]$ and use $d_1(a,b)\le\max(1,2M)d_0(a,b)$ to obtain that $L^1$
error. This gives a PR product modulus. In particular $\mathbf1_A X$
is defined for every represented event $A$ and every represented $X\in L^1$.
This justifies the first-crossing products used in the maximal inequality.

\section{Conditioning, safe event bounds, and continuous transitions}
\label{sec:conditioning}

\subsection{What an upper probability assertion means without an event name}

For a general $L^0$ variable $Y$, the expression $\{Y\ge a\}$ need not
come with a regular event name. A numerical bound on that event should not
be mistaken for a procedure computing its probability. There is nevertheless
an entirely defined observation that expresses the useful estimate.
For rational $0<\eta<a$, put
\[
 r_{a,\eta}(y)
 =\min\left(1,\max\left(0,\frac{y-a+\eta}{\eta}\right)\right).
\]
This bounded Lipschitz function is zero below $a-\eta$ and one at and above
$a$. Its expectation always exists for an $L^0$ name. For $Y\ge0$ in $L^1$,
\[
 \mathbb E r_{a,\eta}(Y)\le\frac{\mathbb EY}{a-\eta}.
\]
The pointwise finite inequality follows by splitting at $a-\eta$ and $a$;
in the transition interval the right side is at least one. For $X\in L^2$,
apply the analogous bound to $r_{a,\eta}(|X|)$ to get
$\mathbb E r_{a,\eta}(|X|)\le\mathbb E X^2/(a-\eta)^2$.

In this report an \emph{upper-envelope conclusion} refers to this family of
defined inequalities with explicit rational slack. For any positive target
slack, a small enough rational $\eta$ is selected from $a$ and the moment
bound by bounded arithmetic. No unspecified outer-measure real is being
postulated. If the hard event has a boundary certificate, the same
observations converge to its indicator with a PR rate and give the usual
probability bound. In standard measure semantics they majorize the hard
event whether or not that event has a computable probability. The same
construction applies to finite maxima, which are continuous functions of
finitely many coordinates.

\subsection{Bayes updating with a certified evidence bound}

Let $\ell\ge0$ be a represented likelihood bounded by $L$, and suppose
$Z=\mathbb E\ell\ge z_0>0$ is certified. On represented bounded
observables define
\[
 \mathbb E_\ell h=\frac{\mathbb E(h\ell)}{Z},\qquad
 \mathbb P_\ell(A)=\frac{\mathbb E(\mathbf1_A\ell)}{Z}.
\]
Bounded multiplication and the reciprocal construction make these regular
real names. Positivity, normalization, and finite additivity follow from the
corresponding expectation identities. Moreover
\[
 |\mathbb P_\ell(A)-\mathbb P_\ell(B)|
 \le (L/z_0)\,\mathbb P(A\triangle B),
\]
so the construction respects completed event equality and transfers
countable-union moduli with the indicated factor.

The numerical stability estimate is explicit. Suppose $|h|\le M$, the
numerator $a=\mathbb E(h\ell)$ is approximated by $a'$ to error
$\delta_a$, and $Z$ by $Z'$ to error $\delta_Z\le z_0/2$. Since
$|a|\le MZ$ and $Z'\ge z_0/2$,
\begin{equation}\label{eq:bayes-budget}
 \left|\frac{a'}{Z'}-\frac aZ\right|
 \le\frac{2(\delta_a+M\delta_Z)}{z_0}.
\end{equation}
This follows by adding and subtracting $a/Z'$ and bounding the two terms.
Conditioning on an event is the case $\ell=\mathbf1_B$. Bayes' identity is
then an equality of intersection probabilities, with every denominator's
positive margin retained. No search for a positive lower bound is inferred
from a bare negation of zero.

For a joint density $f(x,y)$, a pointwise conditional density can be
constructed on a compact observation region if its section integrals have
uniform quadrature and tail data and the marginal
$m(y)=\int f(x,y)\,dx$ has a certified lower bound there. Then
$f(x,y)/m(y)$ is a normalized density in $x$. Finite Fubini, with the
supplied uniform tails, proves the conditional integration identity
\[
 \int_B\left(\int_A\frac{f(x,y)}{m(y)}\,dx\right)m(y)\,dy
   =\int_{A\times B}f(x,y)\,dx\,dy
\]
for the represented integrable tests or regular events at issue. The
Gaussian example below supplies all these data explicitly. General
conditional distributions are not automatically computable from a computable
joint law; counterexamples in the computability literature confirm that
these restrictions are substantive \cite{ackerman}.

\subsection{A finite stopping theorem and a Borel--Cantelli fragment}

Let $M_j$ be an integrable martingale for the first-$j$-bit filtration, with
the conditional identities presented as above. Let $\tau\le N$ be a
stopping rule whose decision $\{\tau>j\}$ is a certified event depending
only on the first $j$ observations. In the simplest internal implementation
these are finite cylinder decisions. On a common finite table,
\[
 M_\tau=M_0+\sum_{j<N}\mathbf1_{\{\tau>j\}}(M_{j+1}-M_j).
\]
Each summand has expectation zero by the conditional identity and past
measurability. Hence $\mathbb E M_\tau=\mathbb E M_0$. Bounded-multiplier
closure justifies the represented products. This is a bounded stopping
theorem; an unbounded stopping time requires additional uniform-integrability
and approximation data before a limit can be taken.

For a jointly PR event family $A_j$, suppose the finite tail sums of its
probabilities are at most $2^{-k}$ after a supplied index $C(k)$.
The admitted tail union $B_k=\bigcup_{j\ge C(k)}A_j$ has mass at most
$2^{-k}$, by finite union bounds and the explicit residual modulus. Every
later $A_j$ is included in $B_k$ in the event-algebra order. These small
tail events are an effective Borel--Cantelli certificate. Its familiar
interpretation is that only finitely many $A_j$ occur almost surely. The
arithmetic assertion is the positive inclusion and small-mass certificate,
not a nonconstructive inference from a logical negation of convergence.

\subsection{Continuous-state Markov chains with supplied transition data}

A useful transition presentation supplies a function $\Psi(x,u)$ taking a
state and an independent uniform input, a finite approximation scheme, and
stability and tail bounds. For example, a coupling estimate
\[
 \mathbb E\min(1,|\Psi(x,U)-\Psi(y,U)|)
 \le C\min(1,|x-y|)
\]
with rational $C\ge1$ propagates $L^0$ approximation errors.
Using independent bit rows, set $X_{j+1}=\Psi(X_j,U_j)$. This constructs
every finite trajectory jointly. If each numerical transition has error
at most $\delta_j$, the finite recurrence
\[
 e_{j+1}\le Ce_j+\delta_j
 \quad\Longrightarrow\quad
 e_n\le C^ne_0+\sum_{j<n}C^{n-1-j}\delta_j
\]
is a quantifier-free induction. Choosing all input errors at most
$2^{-k}/[(n+1)C^n]$ suffices at horizon $n$. This is a concrete PR
construction of a continuous-state process at finite horizons. An arbitrary
family of finite-dimensional laws, without coherent realizers and effective
tightness, is not a substitute for these data.

For an explicit example take $X_0=x_0$ and
\[
 X_{j+1}=aX_j+b+\sigma G_j,
 \qquad |a|\le r<1,\quad\sigma\ge\sigma_0>0,
\]
where $G_j$ are independent standard normals and all parameters have
represented bounds. Expansion and Gaussian convolution show, for $n\ge1$,
that $X_n$ is normal with
\[
 m_n=a^nx_0+b\sum_{j<n}a^j,\qquad
 v_n=\sigma^2\sum_{j<n}a^{2j}.
\]
These approach $m=b/(1-a)$ and $v=\sigma^2/(1-a^2)$ with bounds
\[
 |m_n-m|\le r^n|x_0-m|,\qquad
 |v_n-v|\le\frac{\sigma^2r^{2n}}{1-r^2}.
\]
The identities $am+b=m$ and $a^2v+\sigma^2=v$, together with the
Gaussian transition identity, show that $N(m,v)$ is invariant under one
transition. This identity is valid also when $a=0$ and does not require
deciding whether the represented parameter $a$ vanishes.
Couple the two normal laws as $m_n+\sqrt{v_n}G$ and $m+\sqrt vG$.
Since $v_n,v\ge\sigma_0^2$ and $\mathbb E|G|\le1$,
\[
 \mathbb E|X_n-X_\infty|
 \le r^n|x_0-m|
       +\frac{\sigma^2r^{2n}}{2\sigma_0(1-r^2)}.
\]
Here the displayed coupling compares laws and is not a claim that the
original noisy trajectory converges to a single random variable.
If $r\le1-q$ for rational $q>0$, the elementary bound
$(1-q)^n\le1/(1+nq)$ yields a PR precision schedule. Thus stationary-law
convergence is available without an abstract ergodic theorem. Continuous-time
path constructions and stochastic integration require further modules and
are not claimed in this package.

\section{Products, convolution, and quantitative limit theorems}
\label{sec:limit-theorems}

\subsection{Independence and product integration}

Independence is supplied as a product construction, or by a proved product
factorization on a generating family; it is not inferred from a drawing or
from uncorrelatedness.  For finite rational laws, product weights are
$w_i v_j$ and both Fubini and
$\mathbb E[f(X)g(Y)]=\mathbb E f(X)\,\mathbb E g(Y)$ are identities of
finite sums.  For represented laws, use rational product approximants and
explicit absolute-integrability tails.  The product error estimate
\[
 |uv-u'v'|\leq |u|\,|v-v'|+|v'|\,|u-u'|
\]
and truncation reduce the statement to the finite identity.  A new
integrability condition is needed when the observable is unbounded.

If independent $X,Y$ have represented densities $f,g$, then their sum has
convolution density
\[
 h(z)=\int_{\mathbb R}f(z-y)g(y)\,dy,
\]
provided the density computation has local moduli and uniform tails on the
requested compact $z$-range.  This condition is automatic for the four
explicit families above and for a bounded uniformly continuous $f$ with a
known modulus and an effectively tight $g$: the discarded integral is at
most $\|f\|_\infty$ times the tail mass of $g$.  Normalization and the
pushforward formula are finite double-sum identities followed by the
specified tail limits.  Moreover $h\leq\|f\|_\infty$.

For independent centered normal variables with scales $\sigma,\tau>0$,
completing the square gives
\[
 \frac{(z-y)^2}{\sigma^2}+\frac{y^2}{\tau^2}
 =\frac{z^2}{\sigma^2+\tau^2}
  +\frac{\sigma^2+\tau^2}{\sigma^2\tau^2}
       \left(y-\frac{\tau^2z}{\sigma^2+\tau^2}\right)^2.
\]
The affine substitution in the last Gaussian integral proves
$X+Y\sim N(0,\sigma^2+\tau^2)$.  Every tail needed here is bounded by
\eqref{eq:gaussian-tail}; no characteristic-function uniqueness theorem is
required.  Induction gives
$n^{-1/2}\sum_{i=1}^nG_i\sim N(0,1)$.

\subsection{A Gaussian conditional density}

A useful continuous conditioning example follows from the same algebra.
Let $X,G$ be independent standard normals, let $\sigma\geq\sigma_0>0$, and
put $Y=aX+\sigma G$.  Write $A=a^2+\sigma^2$.  The joint density is
\[
 f(x,y)=\frac1{\sigma Z^2}
  \exp\!\left[-\frac12\left(x^2+\frac{(y-ax)^2}{\sigma^2}\right)\right],
\]
and the marginal of $Y$ is
$m(y)=(\sqrt A Z)^{-1}e^{-y^2/(2A)}$.  Completing the square gives
\[
 x^2+\frac{(y-ax)^2}{\sigma^2}
 =\frac{A}{\sigma^2}\left(x-\frac{ay}{A}\right)^2+\frac{y^2}{A}.
\]
Thus the density kernel $f(x,y)/m(y)$ is a normal density with mean
$ay/A$ and variance $\sigma^2/A$.  It is a defined conditional-density
version at every represented $y$.  It is not formed by dividing by the
zero probability of the singleton event $\{Y=y\}$.

The denominator separation is explicit on compact parameter ranges.
If $A\leq U^2$ with rational $U\geq1$ and $|y|\leq R$, then
\[
 m(y)\geq
 \frac1{4U\,3^{\lceil R^2/(2\sigma_0^2)\rceil}}.
\]
The finite product identity and the same Gaussian tails establish the
conditional integration identity against bounded effectively integrable
test functions.  The lower bound displayed above gives the local precision
required for division and makes the usual Bayesian Gaussian update a
primitive recursive operation on the supplied presentations.

\subsection{Markov, Chebyshev, and the weak law}

For a nonnegative finite observable $Y$ and rational $a>0$,
$a\mathbf1_{\{Y\geq a\}}\leq Y$.  Multiplying by weights and summing gives
Markov's inequality.  For a centered $X$ this yields
\[
 \mathbb P\{|X|\geq\varepsilon\}
 \leq\frac{\mathbb E X^2}{\varepsilon^2}.
\]
For pairwise independent centered $X_1,\ldots,X_n$ with
$\mathbb E X_i^2\leq V$, finite expansion and vanishing cross terms give
$\mathbb E(\sum_iX_i)^2\leq nV$.  Consequently
\begin{equation}
 \mathbb P\!\left\{\left|\frac1n\sum_{i=1}^nX_i\right|
                \geq\varepsilon\right\}
 \leq\frac{V}{n\varepsilon^2}.
 \label{eq:weak-law}
\end{equation}
A requested error probability $\delta>0$ is achieved by
$n\geq\max(1,\lceil V/(\varepsilon^2\delta)\rceil)$.
For noncentered identically distributed variables apply this to
$X_i-\mu$ with a supplied mean and variance certificate.

For general $L^2$ names, a sharp threshold need not itself be an effective
event.  The universal conclusion is then a certified upper bound, obtained
from finite approximants with rational threshold slack and the $L^2$
approximation error.  For the continuous density laws above, sharp threshold
events are available: a bounded density controls boundary strips.  In a sum
of independent density variables, convolution preserves any one available
global density bound; scaling by $1/n$ multiplies that bound by $n$.
Thus every finite threshold or finite maximum event in this section has
an explicit approximation modulus for these laws.  This is precisely what
is needed to pass the finite inequalities to the represented variables.

\subsection{A strong law with an explicit finite-window modulus}

Let $X_i$ be independent, centered, and satisfy
$\mathbb E X_i^2\leq V$.  The input is a single primitive recursive family
of coordinate and finite-product presentations, with the needed moment and
boundary certificates uniform in the coordinate index and finite dimension.
Separate, nonuniform external descriptions of each $X_i$ would not suffice.
Write $S_j=\sum_{i=1}^jX_i$.
For every $m\geq1$ and $a>0$, the finite Kolmogorov maximal inequality is
\begin{equation}
 \mathbb P\{\max_{1\leq j\leq m}|S_j|\geq a\}
 \leq mV/a^2.
 \label{eq:kolmogorov}
\end{equation}
Here full independence, or the appropriate conditional-zero-increment
property, is used; pairwise independence alone does not justify the proof.

\begin{proof}
Let $A_j$ be the event that the threshold is first reached at time $j$.
The events are disjoint, their union is the maximum event, and on $A_j$,
$S_j^2\geq a^2$.  Since $A_j$ and $S_j$ depend only on the first $j$
coordinates, independence and centering give
$\mathbb E[\mathbf1_{A_j}S_j(S_m-S_j)]=0$.  Expansion therefore yields
\[
 \mathbb E[\mathbf1_{A_j}S_m^2]
 =\mathbb E[\mathbf1_{A_j}S_j^2]
  +\mathbb E[\mathbf1_{A_j}(S_m-S_j)^2]
 \geq a^2\mathbb P(A_j).
\]
Sum over $j$ and use $\sum_j\mathbf1_{A_j}\leq1$ and
$\mathbb E S_m^2\leq mV$.  On rational finite spaces this is only an
identity of weighted sums plus positivity.  Moment and boundary certificates
justify the same assertion for the continuous laws.
\end{proof}

For any integers $1\leq N\leq M$, split $[N,M]$ into the dyadic blocks
$[2^jN,2^{j+1}N)$.  Failure of $|S_n|/n<\varepsilon$ in block $j$ implies
$\max_{n\leq2^{j+1}N}|S_n|\geq2^jN\varepsilon$.  By
\eqref{eq:kolmogorov} its probability is at most
$2^{1-j}V/(N\varepsilon^2)$.  The finite geometric sum is less than $4$;
therefore
\begin{equation}
 \mathbb P\!\left\{\max_{N\leq n\leq M}\frac{|S_n|}{n}
                         \geq\varepsilon\right\}
 \leq\frac{4V}{N\varepsilon^2}.
 \label{eq:strong-window}
\end{equation}
In particular, for $\varepsilon=2^{-r}$ and error probability $2^{-k}$,
\begin{equation}
 N(r,k)=\max\!\left(1,\left\lceil4V\,2^{2r+k}\right\rceil\right)
 \label{eq:strong-modulus}
\end{equation}
works \emph{uniformly in every finite endpoint $M\geq N(r,k)$}.
This is a substantive finite, primitive recursive form of the strong law.

When each finite deviation event has a uniform effective event name,
there is also a precise almost-sure certificate in the completed event
algebra.  Put $D_n^\varepsilon=\{|S_n|/n\geq\varepsilon\}$ and let
$F_{N,T}^\varepsilon$ be its union for $N\leq n\leq T$.  For $T'\geq T$,
\eqref{eq:strong-window} bounds the distance between the two finite unions by
\[
 d(F_{N,T'}^\varepsilon,F_{N,T}^\varepsilon)
 \leq\frac{4V}{(T+1)\varepsilon^2}.
\]
Thus their event-metric limit $A_N^\varepsilon$ exists with an explicit
primitive recursive modulus, and
$\mathbb P(A_N^\varepsilon)\leq4V/(N\varepsilon^2)$.  This supplies a
countable union by its residual-event bound; it does not incorrectly sum
the generally divergent individual Chebyshev estimates.

For each error parameter $p$, define the admitted union
\[
 U_p=\bigcup_{r\geq0}A_{N(r,p+r+1)}^{\,2^{-r}}.
\]
The $r$th event has mass at most $2^{-p-r-1}$, so this union has a geometric
residual bound and $\mathbb P(U_p)\leq2^{-p}$.  For every $r$ and every
$n\geq N(r,p+r+1)$, the deviation event $D_n^{2^{-r}}$ is included in
$U_p$, in the order of the event algebra.  This is an explicit almost-sure
convergence certificate: one small exceptional event simultaneously
controls all accuracy scales, with a supplied modulus outside it.
Under an ordinary set realization it implies the usual strong law.
The first-order content is the uniform finite-window estimate, the
explicitly coded exceptional events, and these inclusion estimates; it is
not quantification over an unrepresented powerset of all infinite sample
sequences.  For general $L^2$ names without threshold-event certificates,
retain the upper-envelope finite-window conclusion rather than claiming
these additional event names.

\subsection{An elementary quantitative central limit theorem}

The following rate is intentionally weaker than the optimal classical
Berry--Esseen rate; its advantage is a short proof with completely explicit
primitive recursive bounds.

Let $X_1,\ldots,X_n$ be independent and identically distributed, with
\[
 \mathbb E X_i=0,\qquad \mathbb E X_i^2=1,\qquad
 \mathbb E|X_i|^3\leq B,
\]
where $B$ is a supplied rational bound and the third absolute moment has
an effective tail certificate.  Let $G_i$ be independent standard normal
variables, independent also of all the $X_i$.  If a represented test
function $h$ is twice continuously differentiable, with second derivative
$L$-Lipschitz, then
\begin{equation}
 \left|\mathbb E h\!\left(n^{-1/2}\sum_{i=1}^nX_i\right)
               -\mathbb E h(G)\right|
 \leq\frac{L(B+2)}{6\sqrt n},
 \label{eq:smooth-clt}
\end{equation}
provided the displayed expectations are effectively integrable.  Bounded
$h,h',h''$, as in the application below, suffice.

\begin{proof}
The integral Taylor remainder gives, for all $x,u$,
\[
 h(x+u)=h(x)+h'(x)u+\tfrac12h''(x)u^2+R(x,u),\qquad
 |R(x,u)|\leq L|u|^3/6.
\]
This follows directly by integrating
$|h''(x+t)-h''(x)|\leq L|t|$ twice; the calculation is valid for negative
$u$ after reversing the integration limits.  It can alternatively be
proved by rational Riemann sums with the same error bound.
Replace the summands one at a time.  For replacement $i$, let $T_i$ be the
sum of the other normalized coordinates, using already replaced Gaussian
coordinates and the still unreplaced $X$ coordinates.  It is independent
of both $X_i$ and $G_i$.  Conditional on a fixed value of $T_i$, the constant,
linear, and quadratic Taylor terms have equal expectations for
$u=X_i/\sqrt n$ and $u=G_i/\sqrt n$, since their first two moments match.
The two remainder expectations have total magnitude at most
\[
 \frac{L}{6n^{3/2}}
      \bigl(\mathbb E|X_i|^3+\mathbb E|G_i|^3\bigr)
 \leq\frac{L(B+2)}{6n^{3/2}}.
\]
Product integration justifies averaging this bound in $T_i$.  Sum the $n$
replacement errors and use the proved Gaussian convolution identity.
For a direct arithmetic implementation, first truncate and approximate all
coordinates with the supplied third-moment tails, prove the finite
replacement bound with its recorded moment errors, and allocate the
requested extra precision across the $n$ replacements.  The exact matching
identities are recovered by those explicit error bounds, not by silently
asserting that rounded finite approximants have exactly unit variance.
\end{proof}

To turn this into a distribution-function theorem, use the explicit cutoff
\[
 H(u)=\begin{cases}
 1,&u\leq0,\\
 1-10u^3+15u^4-6u^5,&0\leq u\leq1,\\
 0,&u\geq1.
 \end{cases}
\]
It satisfies $0\leq H\leq1$, because
$H'(u)=-30u^2(1-u)^2\leq0$ in the transition interval.
The function is $C^2$, and $H''$ is globally $60$-Lipschitz: its third
derivative in $(0,1)$ is $-60+360u-360u^2$, which lies between $-60$ and
$30$, while its second derivative vanishes at both endpoints and outside
the interval.  No assertion of a globally continuous third derivative is
needed.

For $h_{t,\eta}(x)=H((x-t)/\eta)$ the Lipschitz constant is $60\eta^{-3}$,
and
\[
 \mathbf1_{(-\infty,t]}\leq h_{t,\eta}
 \leq\mathbf1_{(-\infty,t+\eta]}.
\]
The analogous cutoff starting at $t-\eta$ gives the lower bracket.  Since
$\varphi\leq1$, the Gaussian mass in an interval of length $\eta$ is at most
$\eta$.  Applying \eqref{eq:smooth-clt} to both brackets gives
\begin{equation}
 \sup_{t\in\mathbb R}|F_n(t)-\Phi(t)|
 \leq\eta+\frac{10(B+2)}{\sqrt n\,\eta^3}.
 \label{eq:cdf-clt}
\end{equation}
The supremum notation abbreviates the uniform inequality for every
represented real $t$; no operation computing an arbitrary supremum is
required.  For a requested error $\delta=2^{-k}$, set $\eta=\delta/2$.  It is
enough to take
\begin{equation}
 n\geq N_{\rm CLT}(k)
 :=\left\lceil160(B+2)\right\rceil^2\,2^{8k}.
 \label{eq:clt-rate}
\end{equation}
More generally the smooth theorem holds without identical distributions if
each summand has mean zero, variance one, and the same third-moment bound.

For continuous summand laws with the boundary certificates already given,
$F_n$ is an effective distribution function, so \eqref{eq:cdf-clt} is an
ordinary quantitative CDF assertion.  Without that event information, the
same proof supplies the upper and lower smooth-probability brackets.  It
does not manufacture a CDF name at a threshold whose boundary mass lacks the
required certificate.  For an original
law with mean $\mu$ and variance $\sigma^2>0$, supply a rational lower bound
on $\sigma$ and apply the theorem to $(X_i-\mu)/\sigma$.

\subsection{What remains to be formalized for these results}

The finite arithmetic is not starting from zero.  In the inspected ProveIt
snapshot, the directory
\texttt{Derivations/Lean/IPRADerivations/Probability/} contains object-level
\texttt{PRAAx} derivations for finite probability.  In particular,
\texttt{WeakLaw.lean} contains \texttt{weak\_law},
\texttt{FirstCrossing.lean} contains \texttt{cover}, \texttt{disjoint}, and
\texttt{crossing\_bound}, and the \texttt{Kolmogorov*} files assemble the
finite maximal argument.  These are relevant providers, not merely
mathematical prose.

The limits of that reuse matter.  \texttt{NormalMoments.lean} proves the
arithmetic of double factorials and explicitly says that the integrals
identifying them with normal moments are not formalized.
\texttt{CompleteSquare.lean} proves finite polynomial identities and
explicitly excludes normalization integrals.  \texttt{ProductReplacement.lean}
proves a finite product telescope and bounds, while its header excludes
characteristic functions and smoothing.  The Lindeberg proof here uses a
sum of expectation differences rather than that characteristic-function
product theorem, though the finite telescope infrastructure remains useful.

A completion proceeds in the following concrete order.  First define
rational quadrature and its finite error predicate, then compile the
exponential series estimates and the kernel-tail contracts.  Next construct
normalization and reciprocal names with their separation constants, and
prove moment identities by finite integration-by-parts estimates.  Prove
product integration with explicit truncation errors, boundary-strip transfer,
and Gaussian convolution.  The finite weak and maximal inequalities can
then be transported to the represented variables.  Finally encode the
quintic cutoff, its $60$-Lipschitz second-derivative certificate, the Taylor
remainder estimate, and the one-coordinate replacement inequality.  A
quantifier-free induction on the number of replacements proves the finite
sum bound.  Every precision choice in this section is an explicit primitive
recursive expression or a bounded search with a displayed bound.

\section{Lean definitions, object-language targets, and completion plan}
\label{sec:formalization}

\subsection{What has actually been compiled}

The archive contains a standalone Lean 4.19.0 project using only the bundled
\texttt{Std} library. Its build was run successfully. The two finite
mathematics modules contain 30 named proved theorems, and their build also
checks three closed examples. The other modules contain compiling
primitive-recursive syntax and representation specifications. There are no
\texttt{sorry}, \texttt{admit}, or new axiom declarations in the source.
Representative axiom reports use \texttt{propext} and, for some results,
\texttt{Quot.sound}; none of those inspected reports uses
\texttt{Classical.choice} or \texttt{sorryAx}. The reports are included.

These are semantic Lean proofs of finite algebra. They are not described as
completed object-iPRA derivations of the analytic results. ProveIt's existing
object calculus uses a separately pinned toolchain; its imported sources were
inspected but not rebuilt for this package. The standalone project deliberately
has no hidden dependency on that checkout.

\begin{longtable}{>{\raggedright\arraybackslash}p{.26\textwidth}>{\raggedright\arraybackslash}p{.65\textwidth}}
\toprule
Compiled file & Verified content\\
\midrule\endhead
\texttt{Finite.lean} & Bounded sums, finite event mass, complement and disjoint
additivity, finite union bounds, product rectangles, Markov, a centered
power-Markov form of Chebyshev, and stability of event mass in symmetric
difference.\\
\texttt{L1Finite.lean} & Signed finite sums, the absolute-sum inequality,
expectation stability, and the finite $L^1$ triangle inequality.\\
\texttt{PRSyntax.lean} & Arity-indexed PR syntax, finite evaluation derivations,
and a connection between one fixed syntax tree and its semantic function.
No universal PR evaluator is defined.\\
\texttt{Specifications.lean} & Positive-denominator rational codes, fixed-PR
rational sequences, regular real specifications, PR event tables, common
prefix refinement, and a regular event specification.\\
\bottomrule
\end{longtable}

The compiled regular-real and event specifications use the stronger
one-sided convention $d(x_k,x_l)\le2^{-k}$ for $l\ge k$.
The report's symmetric convention is $d(x_k,x_l)\le2^{-k}+2^{-l}$.
A one-bit shift converts a symmetric name to the one-sided convention, and
the latter already implies the symmetric one. This harmless reindexing is
explicit and should be retained when connecting the interfaces.

\subsection{All representation data for the continuous layer}

The following additional structures are specifications to implement; they
are not postulated instances. Their arithmetic contents are completely
determined by the definitions in the report.

\begin{longtable}{>{\raggedright\arraybackslash}p{.22\textwidth}>{\raggedright\arraybackslash}p{.69\textwidth}}
\toprule
Interface & Fields and exact obligation\\
\midrule\endhead
\texttt{PRTable} & A fixed PR depth program $\ell(k)$ and three fixed PR
programs on $(k,i)$ for positive/negative numerator and denominator minus
one. Only $i<2^{\ell(k)}$ is read. Prefix lifting uses
$i\mapsto\lfloor i/2^{D-\ell(k)}\rfloor$.\\
\texttt{RandomName0} & A \texttt{PRTable} and a derivation of the all-precision
$d_0$ inequality, where $d_0$ is the finite rational fold of
$\min(1,|x-y|)$.\\
\texttt{RandomName1} & The same, using the fold of $|x-y|$. Expectation is
the rational sum divided by the row count.\\
\texttt{RandomName2} & The same, with
$I((x_k-x_l)^2)\le(2^{-k}+2^{-l})^2$. No square-root test is needed
in the validity predicate.\\
\texttt{EventName} & Boolean tables and symmetric-difference counts;
the cleared formula is given in Section~\ref{sec:scope}.\\
\texttt{BandSpec} & Fixed PR $b(k),R(k)$ and proofs that for every
$r\ge R(k)$ the finite band count at width $2^{-b(k)}$ is at most $2^{-k}$.\\
\texttt{CompactFunction} & A fixed rational evaluator, local bound, continuity
modulus, and the evaluator's finite consistency and oscillation proofs.
Optional increment and derivative contracts support calculus.\\
\texttt{KernelSpec} & A nonnegative \texttt{CompactFunction} on each compact
range, its listed breakpoints, mass-tail radius, and a fixed positive
normalization witness (K1)--(K2).\\
\texttt{MomentSpec} & An integer order $p$, PR cutoff, and finite annular-tail
proofs for $|x|^p w(x)$. A numerical bound on the moment is additional data
when a theorem, such as the CLT, requests one.\\
\texttt{ProductFamily} & A fixed PR bit-position map, injectivity/disjoint-row
counting proofs, and uniform finite-dimensional approximation schedules.
These implement the independent copies used here.\\
\texttt{SmoothTest} & PR evaluators for $h,h',h''$, bounds, and a proved
Taylor remainder $|R(x,u)|\le L|u|^3/6$, with the finite evaluation errors
included. The quintic cutoff supplies a concrete instance.\\
\texttt{TailUnionSpec} & A jointly PR event family and a PR rate for distances
between its finite unions. A summable mass-tail provider is one constructor.\\
\bottomrule
\end{longtable}

For clarity, the compiled syntax begins with the following shape (the source
contains its fully typed Lean form):
\begin{lstlisting}
inductive PR : Nat -> Type where
  | zero (n : Nat) : PR n
  | succ : PR 1
  | proj (n : Nat) (i : Fin n) : PR n
  | comp (outer : PR m) (inner : Fin m -> PR n) : PR n
  | primrec (base : PR n) (step : PR (n + 2)) : PR (n + 1)
\end{lstlisting}
A finite \texttt{Eval} derivation has separate constructors for the defining
equations. \texttt{FixedPRFunction} carries one code, its semantic function,
and evaluation-graph proofs for that fixed code. The function field is a
metatheoretic convenience, not an admission of all Lean functions into iPRA.
The final object version must carry quoted PR symbols and actual proofs of
the numerical contracts in \texttt{Prf PRAAx}.

\subsection{The main embedded theorem targets}

For a fixed table program $X$ let $\operatorname{Reg}_0(X)$ denote the
all-precision first-order formula above, and similarly for the other
metrics. Let $\operatorname{Band}(X,t)$ abbreviate the quantified finite
band test with its fixed $b,R$ programs. The intended targets are genuine
formulas over natural codes, such as
\begin{align*}
 &\operatorname{Reg}_1(X)\ \longrightarrow\
   \operatorname{Reg}_{\mathbb R}(\operatorname{Expect}(X)),\\
 &\operatorname{Reg}_0(X)\wedge\operatorname{Band}(X,t)
   \ \longrightarrow\
   \operatorname{Reg}_{\rm ev}(\operatorname{Threshold}(X,t)),\\
 &\operatorname{Kernel}(w)\ \longrightarrow\
   \operatorname{Reg}_0(\operatorname{HistogramSampler}(w)),\\
 &\operatorname{Kernel}(w)\wedge\operatorname{Moment}_2(w)
   \ \longrightarrow\
   \operatorname{Reg}_2(\operatorname{ReindexSampler}(w)).
\end{align*}
Each named construction on the right is a fixed PR program built from the
specified input definitions. There is no free functional quantifier in these
schemas. Event and random-variable equality are all-precision relations;
they are not tested by an evaluator.

For the strong law, quote the finite sum, maximum, event indicator, product
weight, and error-allowance programs. The endpoint is the cleared version of
\eqref{eq:strong-window}, universally quantified over $N,M$ and rational
precision parameters under the fixed presentation contracts. Its derived
union name uses $T\ge4V\,2^k/\varepsilon^2$ for residual precision
$2^{-k}$. For the CLT, quote the quintic cutoff and the finite Taylor-error
test; the target is a pair of expectation inequalities whose error is
\eqref{eq:cdf-clt}, or its squared version with all positive square roots
eliminated by rational intervals. The sample-size condition
\eqref{eq:clt-rate} is already integral arithmetic.

In ProveIt the final endpoint has the form
\begin{lstlisting}
-- Desired proof type after the stated formula has been constructed:
Prf PRAAx ([] : List (Fm PTm 0)) phi
\end{lstlisting}
This is a target type, not an asserted theorem constant in the delivered
code. A law field in a structure, an axiom, a standard-model truth proof,
or a program that merely checks examples would not fill this endpoint.

\subsection{A finite dependency plan with explicit induction invariants}

\begin{longtable}{>{\raggedright\arraybackslash}p{.08\textwidth}>{\raggedright\arraybackslash}p{.38\textwidth}>{\raggedright\arraybackslash}p{.45\textwidth}}
\toprule
Stage & Construction and proof & Quantifier-free invariant or finite bridge\\
\midrule\endhead
1 & Quote rational codes, list access, refinement, sums and norms. &
Valid denominators, running sum identity/inequality, and prefix projection
composition. Reuse the existing signed and bounded-sum arithmetic.\\
2 & Prove finite probability and metric estimates. &
Weighted partial-sum bounds; rowwise indicator identities; expanded
nonnegative squares. Existing \texttt{Expectation} providers already
cover much of this layer.\\
3 & Reify regularity and extensional equality. &
Instantiate input contracts at requested finite indices, use triangle
inequalities, then remove rational slack. Do not induct on a quantified
regularity predicate.\\
4 & Quote compact quadrature and calculus. &
Accumulated oscillation/error budget; telescoping increment identity;
partial Taylor sum with its geometric tail.\\
5 & Normalize kernels and realize laws. &
Nonnegative dyadic masses, sum at most one, residual at zero, monotone
cumulative allocation, finite clipped coupling and moment-tail bounds.\\
6 & Extend events, products and conditioning. &
Symmetric-difference budget, bounded-multiplier truncation bound, finite
double-sum exchange, positive denominator guard.\\
7 & Lift weak and maximal inequalities. &
First-crossing partition and centered product identities; squared error
and boundary-band budgets before the limit. Reuse \texttt{WeakLaw},
\texttt{FirstCrossing}, and the explicitly scoped \texttt{Kolmogorov*} providers.\\
8 & Construct strong-law exceptional events. &
Finite dyadic-block geometric sum, finite-union residual estimate,
then explicit diagonal schedule for the countable precision scales.\\
9 & Prove the quantitative CLT. &
After $j$ replacements, the expectation discrepancy is bounded by the
sum of the first $j$ Taylor/moment errors. Verify the polynomial cutoff
identities and allocate the finite moment-rounding errors.\\
10 & Export and audit. &
Construct the actual closed \texttt{Prf PRAAx} endpoints, then use the
existing transport theorem to the chosen concise presentation. Check
axiom and dependency reports at those endpoints.\\
\bottomrule
\end{longtable}

Stages 1--2 have substantial existing object-level support, plus the new
standalone finite checks. The new continuous constructions in stages 3--9
are supplied here as human proofs and formalization specifications. Their
full embedded implementation remains work to do. The finite examples and
compiled specification records do not turn that remaining work into a
completed formal proof.

\clearpage
\section{Coverage, evidence, and explicit limits}
\label{sec:coverage}

\begin{longtable}{>{\raggedright\arraybackslash}p{.32\textwidth}>{\raggedright\arraybackslash}p{.58\textwidth}}
\toprule
Topic & Result and exact scope\\
\midrule\endhead
Probability space and events & Coherent rational cylinders, represented
measure algebra, finite probability axioms, countable additivity for
supplied finite-union tail moduli.\\
Random variables and expectation & $L^0$, $L^1$, $L^2$ names; law through
bounded observables; expectation, covariance, product and multiplier
closure under their stated integrability assumptions.\\
Continuous distributions & Uniform, exponential, normal, Cauchy, and Cantor
singular continuous law; explicit tails, moments when finite, CDFs,
histogram realization, and separated pointwise quantiles.\\
Joint laws and conditioning & Explicit independent copies, represented
Fubini with absolute tails, convolution, finite conditional averaging,
Bayes updates with positive evidence, Gaussian conditional density.\\
Convergence & PR-modulated completion, quantitative Vitali and dominated
convergence, canonical dyadic-martingale convergence, certified first
Borel--Cantelli fragment.\\
Laws of large numbers & Finite weak-law bound, finite-window maximal/strong
bound, and effective exceptional-event certificates when threshold events
have the required presentation.\\
Central limit theorem & Smooth Lindeberg replacement and uniform CDF error
with explicit $O(\delta^{-8})$ sufficient sample size under a third-moment
bound. An optimal Berry--Esseen estimate is not claimed.\\
Stochastic processes & Finite-horizon continuous-state transitions, bounded
stopping, and an explicit autoregressive Gaussian stationary-law example.
No general continuous-time path or stochastic-integral construction.\\
\bottomrule
\end{longtable}

The Python reference uses integer and rational arithmetic in all acceptance
checks. The recorded run passes 2,963 checks: cylinder refinements and
moments, conditional averaging and tower identities, independence,
clipping, finite first-crossing inequalities, cutoff-polynomial identities,
and certified Gaussian enclosures. In particular it certifies
\[
 \frac{882141}{1048576}\le\Phi(1)\le\frac{55143}{65536},
 \qquad 0.841<\Phi(1)<0.842.
\]
The decimal endpoints on the right abbreviate exact rationals. The narrower
fraction interval is obtained by positive Taylor sums, inversion, midpoint
quadrature, and the Gaussian tail inequality. Floating-point libraries are
not proof oracles in this calculation. Such finite checks validate the
implementation and constants; the general conclusions rest on the human
proofs and, for the finite Lean layer, kernel checking.

For perspective, $V=1$, $\varepsilon=1/10$, $\delta=1/20$ give sufficient
sample sizes $2{,}000$ for the weak law and $8{,}000$ for the uniform
strong-law window. The elementary CLT with $B=2$ and desired error $1/10$
gives the much larger sufficient value $40{,}960{,}000{,}000{,}000$.
That conservative constant is a cost of the elementary smoothing proof,
not a statement about the actual convergence speed of a particular law.
The package's purpose is a transparent proof-theoretic foundation; sharper
certified algorithms can replace these providers without changing the
representation interfaces.

The precise exclusions are as informative as the positive results.
Arbitrary measurable-set comprehension, regular probabilities for all
effective Borel descriptions, general rate extraction from qualitative
convergence, and unrestricted disintegration are not available for free.
Their absence is supported by the explicit halting obstructions and the
conditioning results cited above, rather than by a claim that continuous
mathematics itself is impossible in iPRA. The continuous mathematics succeeds
when the finite information consumed by each operation is stated and proved.

\begin{thebibliography}{9}
\addcontentsline{toc}{section}{References and source provenance}

\bibitem{wikiProbability}
Wikipedia contributors, \emph{Probability theory}, requested topic outline,
accessed 20 November 2026.
\url{https://en.wikipedia.org/wiki/Probability_theory}.
Used to choose coverage, not as a proof oracle.

\bibitem{chan}
Yuen-Kwok Chan, \emph{Foundations of Constructive Probability Theory},
arXiv:1906.01803v2, 2019.
\url{https://arxiv.org/abs/1906.01803}.
Background on Bishop-style constructive probability; no blanket reduction
of that book to iPRA is assumed here.

\bibitem{ackerman}
Nathanael L. Ackerman, Cameron E. Freer, and Daniel M. Roy,
\emph{On the computability of conditional probability},
Journal of the ACM \textbf{66}(3), article 23, 2019, pp.\ 23:1--23:40.
\url{https://arxiv.org/abs/1005.3014};
\url{https://doi.org/10.1145/3321699}.

\bibitem{repoCore}
Vladimir Reshetnikov and ProveIt contributors,
\emph{Intuitionistic primitive recursive arithmetic: ConciseIPRA and
IPRADerivations}, repository source inspected at commit
\texttt{b2341af5370cd8a43f72ebb2341c0d7d7395b634}.
\href{https://github.com/VladimirReshetnikov/ProveIt/tree/b2341af5370cd8a43f72ebb2341c0d7d7395b634/Logic/PrimitiveRecursiveArithmetic}{Pinned iPRA directory}.
The exact selected files, Git blob hashes, and SHA-256 hashes are in
\texttt{metadata/source\_manifest.json}. Object-level probability providers
and their limits are mapped in \texttt{metadata/PROVIDERS.md}.

\bibitem{repoProbability}
ProveIt research packets,
\emph{Expected value in iPRA} (16 November 2026),
\emph{Normal distribution in iPRA} (16 November 2026), and
\emph{Constructive central limit theorems in iPRA} (15 November 2026).
\href{https://github.com/VladimirReshetnikov/ProveIt/tree/b2341af5370cd8a43f72ebb2341c0d7d7395b634/Logic/PrimitiveRecursiveArithmetic/docs/proofs/probability}{Pinned probability packets}.
Their research statements are distinct from the later finite object-language
derivations. The stronger CLT constants in the earlier packet are not used
as premises of this report.

\bibitem{repoObservable}
ProveIt research packet,
\emph{Continuous mathematics in iPRA: shared error budgets, constructive
integration, and stable homotopies}, 20 November 2026, especially
\texttt{foundations.tex}, \texttt{measure.tex}, and
\texttt{model\_probability.tex}.
\href{https://github.com/VladimirReshetnikov/ProveIt/tree/83421593771f877bc1ae98f4229eb35507bff05b/Logic/PrimitiveRecursiveArithmetic/docs/proofs/foundations/ipra_observable_analysis_2026-11-20}{Pinned source packet}.
Its coherent measured algebra, completion interfaces, and represented
integration are reused with attribution.

\bibitem{repoMeasureTransport}
ProveIt research packet,
\emph{Quantitative integration and transport in iPRA},
20 November 2026, especially \texttt{paper/sections/03\_measure.tex}.
\href{https://github.com/VladimirReshetnikov/ProveIt/tree/83421593771f877bc1ae98f4229eb35507bff05b/Logic/PrimitiveRecursiveArithmetic/docs/proofs/foundations/ipra_measure_transport_2026-11-20}{Pinned source packet}.
The corrected clipping estimate, effective tail interfaces, and separation
of finite Lean proofs from embedded analytic derivations are reused.

\end{thebibliography}

\end{document}
